\documentclass[11pt,a4paper]{article}

% ── Packages ─────────────────────────────────────────────────────────────────
\usepackage[utf8]{inputenc}
\usepackage[T1]{fontenc}
\usepackage{lmodern}
\usepackage[margin=2.5cm]{geometry}
\usepackage{amsmath,amssymb,amsthm}
\usepackage{booktabs}
\usepackage{listings}
\usepackage{xcolor}
\usepackage{hyperref}
\usepackage{url}
\usepackage{graphicx}
\usepackage{microtype}
\usepackage{parskip}
\usepackage{authblk}
\usepackage{abstract}

% ── Theorem environments ──────────────────────────────────────────────────────
\newtheorem{theorem}{Theorem}[section]
\newtheorem{lemma}[theorem]{Lemma}
\newtheorem{proposition}[theorem]{Proposition}
\newtheorem{definition}[theorem]{Definition}
\newtheorem{remark}[theorem]{Remark}
\newtheorem{corollary}[theorem]{Corollary}
\newtheorem{assumption}[theorem]{Assumption}

% ── Code listing style ────────────────────────────────────────────────────────
\lstset{
  basicstyle=\ttfamily\small,
  breaklines=true,
  frame=single,
  backgroundcolor=\color{gray!5},
  keywordstyle=\bfseries,
  numbers=left,
  numberstyle=\tiny\color{gray},
}

% ── Metadata ──────────────────────────────────────────────────────────────────
\title{\textbf{NAPQES: A Single-Primitive Authenticated Encryption Scheme\\
Built from HMAC-SHA256}}

\author[1]{Abdeslam Jakjoud}
\author[1]{Karim Amor}
\affil[1]{QAegis, France}
\date{July 2026}

% ─────────────────────────────────────────────────────────────────────────────
\begin{document}
\maketitle

% ── Abstract ─────────────────────────────────────────────────────────────────
\begin{abstract}
We present NAPQES, an Authenticated Encryption with Associated Data (AEAD)
scheme whose security reduces exclusively to the pseudorandomness of
HMAC-SHA256.  Unlike AES-GCM and ChaCha20-Poly1305, NAPQES employs no
finite-field arithmetic, group operations, or AES S-box permutations, and
therefore admits no analogue of the algebraic key-recovery attacks that
follow from nonce reuse in a polynomial authenticator. Like every
symmetric AEAD scheme, and like the two above, it is unaffected by Shor's
algorithm; its post-quantum posture is governed by Grover search alone
(Section~\ref{sec:pq}).  The construction keys two independently sampled
secrets: an
arithmetic-layer key $\mathbf{k}$, an ordered tuple of primes driving a
prime-indexed token cipher, and a uniform 256-bit secret $\mathsf{sk}$ that
keys every HMAC domain derivation used for noise injection, plaintext
padding, keystream masking, and message authentication. The nonce is
derived synthetically as
$N=\mathrm{HMAC}(\mathsf{sk},\texttt{0x0A}\|\mathrm{be8}(|A|)\|A\|M)[0{:}16]$,
in
the manner of AES-GCM-SIV, rather than sampled independently at random.
Because $\mathbf{k}$ and $\mathsf{sk}$ are sampled independently and only
$\mathsf{sk}$ ever keys an HMAC call, we prove IND-CPA-det (the
restricted-query notion appropriate to a deterministic, misuse-resistant
scheme's synthetic nonce), INT-CTXT, and IND-CCA-det security
(Theorems~\ref{thm:ind-cpa}, \ref{thm:int-ctxt}, \ref{thm:ind-cca}) under
the standard, uniform-key PRF assumption on HMAC-SHA256 — with no
key-entropy term and no non-standard key-distribution assumption — and show
that nonce misuse (re-encrypting an identical associated-data/message pair)
degrades gracefully to disclosed message-equality leakage rather than key
recovery. Ciphertext length is proved to be a deterministic function of the
padded-length bucket alone (Lemma~\ref{lem:length-det}), never of message
content, nonce, or associated data. Against a length-observing adversary we
define and prove LH-IND-CPA-det (Theorem~\ref{thm:lh-ind-cpa}), under which
the challenge messages may differ in length, at the same bound; we bound
the residual length leakage at $\log_2 13\approx3.70$ bits per message, and
at zero under the $\mathsf{frame}$ padding profile
(Corollary~\ref{cor:length-leak}); and we exhibit a zero-query adversary
breaking the same notion outright for AES-GCM and ChaCha20-Poly1305. We
further show that this property derives from the padding profile alone and
not from the arithmetic layer, whose removal changes no theorem in the
paper (Propositions~\ref{prop:expansion-neutral}
and~\ref{prop:lean-theorems}). We discuss the post-quantum security
level and implementation notes across
Python, Rust, and C reference implementations, and report the
cross-implementation Known-Answer Test corpora, a NIST SP 800-22 Rev.\ 1a
statistical analysis of ciphertext output, and fuzz coverage of the decoder
— evidence for implementation agreement and decoder robustness
respectively, not for any security claim, which rests on the theorems
alone. A reference implementation, KAT corpus, and fuzz
harnesses are published at \url{https://github.com/Synergeon-epineon/NAPQES}.
\end{abstract}

\tableofcontents
\clearpage

% ─────────────────────────────────────────────────────────────────────────────
\section{Introduction}
\label{sec:intro}
% ─────────────────────────────────────────────────────────────────────────────

Modern AEAD standards—AES-GCM~\cite{nist-gcm} and
ChaCha20-Poly1305~\cite{rfc8439}—have achieved widespread deployment.
Neither is threatened by Shor's algorithm, which applies to hidden-subgroup
problems rather than to symmetric primitives~\cite{nistpq2022}. Their
algebraic structure matters for a different reason: the polynomial
authenticators of both (GHASH over $\mathrm{GF}(2^{128})$, Poly1305) turn a
single nonce reuse into recovery of the authentication key and hence into
arbitrary forgery~\cite{joux2006}. That failure mode, not quantum
cryptanalysis, is what NAPQES's design avoids.

NAPQES takes a more conservative approach: the \emph{only} cryptographic
primitive in the construction is HMAC-SHA256, approved under NIST FIPS
198-1 and FIPS 180-4. The security of the entire scheme therefore reduces
to the pseudorandomness of HMAC-SHA256—a well-studied assumption with no
known quantum speedup beyond Grover's $O(\sqrt{N})$ search, which reduces
the effective security level from 256~bits to approximately 128~bits while
remaining far above all recommended thresholds.

\paragraph{Contributions.}
\begin{enumerate}
  \item We fully specify NAPQES (Section~\ref{sec:construction}): an
        independent HMAC-domain-derivation secret $\mathsf{sk}$ and
        arithmetic-layer key $\mathbf{k}$, sampled independently, together
        with a synthetic, message-bound nonce
        $N=\mathrm{HMAC}(\mathsf{sk},\texttt{0x0A}\|\mathrm{be8}(|A|)\|A\|M)$
        (Definition~\ref{def:synthnonce}) and a fixed-width, per-bucket-ceiling
        wire format (Section~\ref{sec:wire-format}) under which ciphertext
        length depends only on the padded-length bucket.
  \item We give a formal AEAD algorithm triple
        $(\mathsf{KeyGen},\mathsf{Enc},\mathsf{Dec})$
        (Definition~\ref{def:aead-triple}) and prove correctness
        (Proposition~\ref{prop:correctness}).
  \item We prove IND-CPA-det (Definition~\ref{def:ind-cpa-det},
        Theorem~\ref{thm:ind-cpa}), INT-CTXT
        (Definition~\ref{def:int-ctxt}, Theorem~\ref{thm:int-ctxt}), and
        IND-CCA-det (Definition~\ref{def:ind-cca-det},
        Theorem~\ref{thm:ind-cca}) security under the \emph{standard,
        uniform-key}
        PRF assumption on HMAC-SHA256, with no key-entropy term: because
        $\mathsf{sk}$ is sampled uniformly and independently of
        $\mathbf{k}$, every reduction below may sample $\mathbf{k}$ on its
        own to run the arithmetic layer while forwarding every HMAC call to
        an external PRF oracle keyed by $\mathsf{sk}$, with no simulation
        gap of any kind.
  \item We prove ciphertext length is a deterministic function of the
        padded-length bucket alone (Lemma~\ref{lem:length-det}) — never of
        message content within a bucket, the nonce, or the associated data.
  \item We define LH-IND-CPA-det (Definition~\ref{def:lh-ind-cpa-det}), the
        length-hiding strengthening of IND-CPA-det in which the challenge
        messages need only share a padding bucket, prove that NAPQES
        satisfies it at the same bound (Theorem~\ref{thm:lh-ind-cpa}),
        quantify the residual leakage over an arbitrary padding profile
        (Corollary~\ref{cor:length-leak}), and exhibit the zero-query
        adversary that achieves maximal advantage against AES-GCM and
        ChaCha20-Poly1305 in the same game
        (Proposition~\ref{prop:lh-separation}).
  \item We attribute that property to its actual source. The expansion
        factor is a public injective function of the padding bucket and so
        contributes no length confidentiality
        (Proposition~\ref{prop:expansion-neutral}); the arithmetic layer
        contributes to no theorem in this paper, and we specify the same
        construction with that layer removed, together with the bounds it
        preserves and the bandwidth it recovers
        (Section~\ref{sec:lean-profile},
        Proposition~\ref{prop:lean-theorems}).
  \item We provide a post-quantum analysis (Section~\ref{sec:pq}) and a
        comparison with AES-GCM, ChaCha20-Poly1305, and Ascon
        (Section~\ref{sec:comparison}).
  \item We report NIST SP 800-22 Rev.\ 1a statistical test results on
        ciphertext output (Section~\ref{sec:sts}), the cross-implementation
        KAT corpora
        (Section~\ref{sec:kats}), and a fuzz harness covering the full
        decoder pipeline (Section~\ref{sec:fuzz}), and state in each case
        what the evidence does and does not establish.
  \item We discuss known caveats (Section~\ref{sec:caveats}): the 16-bit
        plaintext length cap, the fixed ciphertext expansion factor, the
        padding length-bucket leak, and the scope of what each of
        $\mathsf{sk}$ and $\mathbf{k}$ actually protects.
\end{enumerate}

% ─────────────────────────────────────────────────────────────────────────────
\section{Related Work}
\label{sec:related}
% ─────────────────────────────────────────────────────────────────────────────

\paragraph{AES-GCM~\cite{nist-gcm}.}
AES-GCM combines AES-CTR for confidentiality with a GHASH polynomial MAC
over $\mathrm{GF}(2^{128})$. Under nonce reuse the two keystreams cancel
and GHASH degenerates into a solvable polynomial system, so the
authentication key $H$ is recoverable and arbitrary forgeries follow: this
is the Forbidden Attack~\cite{joux2006}. Separately, Iwata, Ohashi and
Minematsu~\cite{iwata2012} showed that GCM's original security proof was
flawed for non-96-bit nonces and repaired the bounds; that work concerns
the proof, not a nonce-reuse break, and we cite it only for that. NAPQES
uses HMAC-SHA256 for authentication, which does not admit the Forbidden
Attack, and derives its nonce synthetically so that reuse across distinct
$(A,M)$ pairs is not a caller-controlled failure mode.

\paragraph{ChaCha20-Poly1305~\cite{rfc8439}.}
ChaCha20 is an ARX cipher whose security is not reduced to a well-studied
hardness assumption; Poly1305 is again a polynomial MAC, subject to the
same class of Forbidden Attacks as GHASH when the one-time key is reused.
Neither ChaCha20 nor Poly1305 is approved under current FIPS standards.

\paragraph{Ascon~\cite{ascon2021}.}
Ascon (NIST LWC winner) is a sponge-based AEAD targeting lightweight
environments. Its security is based on the hardness of distinguishing the
Ascon permutation from a random permutation, which involves xor-and-rotate
operations over 64-bit words. Ascon has since been standardised by NIST as
SP 800-232~\cite{sp800-232}, so it is an approved algorithm; the
comparison in Section~\ref{sec:comparison} reflects that status.

\paragraph{HMAC-based constructions.}
HKDF~\cite{rfc5869} and HMAC-DRBG~\cite{nist-drbg} demonstrate that
HMAC-SHA256 is a sufficient basis for key derivation and pseudorandom
generation. NAPQES extends this philosophy to the full AEAD context.

\paragraph{Deterministic and misuse-resistant AE.}
Rogaway and Shrimpton~\cite{rogaway2006} formalise deterministic
authenticated encryption (DAE) and the corresponding restricted-query
security notion, realised by SIV~\cite{rfc5297} and
AES-GCM-SIV~\cite{rfc8452}. NAPQES adopts the same synthetic-IV philosophy
for its nonce. It does \emph{not} inherit those schemes' security notion:
what is proved here is the weaker pair of statements of
Section~\ref{sec:security}, as Remark~\ref{rem:dae-scope} sets out.

% ─────────────────────────────────────────────────────────────────────────────
\section{Construction}
\label{sec:construction}
% ─────────────────────────────────────────────────────────────────────────────

\subsection{Notation}
\label{sec:notation}

Throughout this paper:
\begin{itemize}
  \item $\|$ denotes byte concatenation.
  \item $\mathrm{HMAC}(K, M)$ denotes HMAC-SHA256 keyed with $K$ over message $M$.
  \item $\mathrm{be4}(x)$, $\mathrm{be5}(x)$, $\mathrm{be8}(x)$ denote the
        4-, 5-, and 8-byte big-endian encodings of the non-negative integer
        $x$ (each field is fixed-width; $x$ is assumed to fit).
  \item $[0{:}n]_{\mathrm{be}k}$ denotes the first $n$ bytes of a 32-byte
        HMAC-SHA256 digest, interpreted as an unsigned big-endian integer
        in $[0,2^{k})$ where $k=8n$.
  \item $\mathcal{P}$ denotes the set of prime integers in $[10^6, 9.9 \times 10^6]$,
        a fixed public system parameter with $|\mathcal{P}| = 579{,}947$
        (verified by sieve; the interval is inclusive at both ends, and both
        endpoints are composite). Implementations MUST use this interval:
        a different upper bound yields a different $\mathcal{P}$, a
        different canonical enumeration, and hence different keys.
\end{itemize}

\subsection{Key Generation}
\label{sec:keygen}

A NAPQES key pair is $(\mathbf{k}, \mathsf{sk})$, sampled by
$\mathsf{KeyGen}(1^\lambda)$ as follows.

\begin{definition}[Key generation]\label{def:keygen}
Draw $\mathbf{k} = (k_1,\ldots,k_K)$, an ordered tuple of $K$ distinct
primes sampled uniformly, without replacement, from $\mathcal{P}$ via a
cryptographically secure DRBG (default $K=10$); \emph{independently}, draw
$\mathsf{sk}\xleftarrow{\$}\{0,1\}^{256}$ uniformly via a CSPRNG. No
function relates $\mathsf{sk}$ to $\mathbf{k}$. The serialised arithmetic
key is
\[
  \mathsf{key\_bytes}(\mathbf{k}) = \mathrm{be5}(k_1) \| \mathrm{be5}(k_2) \| \cdots \| \mathrm{be5}(k_K).
\]
\end{definition}

\begin{remark}[Concrete sampling algorithm]\label{rem:sampling}
$\mathcal{P}$ is fixed as a public, canonical ascending enumeration
$\mathcal{P}=\{p_0<p_1<\cdots<p_{|\mathcal{P}|-1}\}$ (e.g.\ a precomputed
sieve table), so ``the $idx$-th prime of $\mathcal{P}$'' is well-defined
and identical across implementations. Index sampling is by rejection:
$\textsc{SampleIndex}()$ draws a uniform 32-bit unsigned integer $x$ from
the CSPRNG, sets $L=\lfloor 2^{32}/|\mathcal{P}|\rfloor\cdot|\mathcal{P}|$,
and rejects and redraws $x$ whenever $x\geq L$; once $x<L$, it returns
$idx=x\bmod|\mathcal{P}|$, which is then exactly uniform over
$\{0,\ldots,|\mathcal{P}|-1\}$ with no modulo bias. $k_1$ is set to
$p_{\textsc{SampleIndex}()}$; for $j=2,\ldots,K$, $k_j$ is drawn by
repeatedly invoking $\textsc{SampleIndex}()$ and rejecting (redrawing) any
result equal to an already-chosen $k_1,\ldots,k_{j-1}$, until a fresh
value is obtained. This realises exactly the uniform distribution over
ordered tuples of $K$ distinct elements of $\mathcal{P}$ assumed by the
$H_\infty(\mathbf{k})$ computation below.
\end{remark}

The min-entropy of $\mathbf{k}$ as sampled above is
\[
  H_\infty(\mathbf{k}) = \log_2\!\left(\frac{|\mathcal{P}|!}{(|\mathcal{P}|-K)!}\right) \approx K\log_2|\mathcal{P}| \approx 19.1456\,K \text{ bits},
\]
so $H_\infty(\mathbf{k}) = 191.46$ bits at the default $K=10$
(key-space size $\approx 4.30\times10^{57}\approx 2^{191.46}$). All figures
here are computed from the exact $|\mathcal{P}|=579{,}947$, not from an
approximation to it.

\begin{remark}[The role of $\mathbf{k}$ and $\mathsf{sk}$ are not symmetric]
\label{rem:key-roles}
$\mathsf{sk}$ is the \emph{only} secret that ever keys an HMAC call in this
construction (Section~\ref{sec:domains}); $\mathbf{k}$ is used exclusively,
directly, and un-keyed by the arithmetic token map $c \mapsto c\cdot k_j+a$
(Section~\ref{sec:token-construction}). Consequently, none of the security
theorems of Section~\ref{sec:security} depend in any way on
$H_\infty(\mathbf{k})$: the PRF assumption below is stated, and used,
against the standard \emph{uniform}-key distribution of $\mathsf{sk}$ only.
$K$ therefore carries no security-theorem-driven minimum in this
construction (contrast a single-secret design in which the arithmetic key
also keys HMAC, which would require a minimum $K$ to keep an offline
key-guessing term negligible). We state plainly what $K$ does \emph{not}
do. $K$ is never transmitted, appears nowhere in the wire format, and does
not influence $|C|$, which by Lemma~\ref{lem:length-det} is a function of
the padding bucket alone; in particular $K$ makes no contribution to the
length-hiding property of Section~\ref{sec:length-hiding}, and no claim of
traffic-analysis resistance anywhere in this paper rests on $K$, on
$H_\infty(\mathbf{k})$, or on the multi-prime structure
(Remark~\ref{rem:coarsening-not-inflation}). $K=10$ is an interoperability
default, not a security parameter: any $K\ge1$ is permitted by the
construction, and both parties must agree on it out of band exactly as they
agree on $\mathbf{k}$ itself. The reference implementations emit a warning
below $K=7$; that threshold is inherited from the earlier single-secret
schedule, in which $H_\infty(\mathbf{k})$ did enter the bound, and no
theorem in this paper requires it. What the arithmetic layer is retained
for instead is analysed in Section~\ref{sec:arith-layer-value}.
Section~\ref{sec:caveats} discusses,
conversely, what protecting $\mathbf{k}$ actually buys once $\mathsf{sk}$
is known.
\end{remark}

\subsection{Domain-Separated HMAC Derivation}
\label{sec:domains}

All internal values are computed as
\[
  \mathrm{Derive}_d(\mathsf{sk}, N, \mathsf{ctx}) = \mathrm{HMAC}(\mathsf{sk},\ d \| N \| \mathsf{ctx}),
\]
where $d$ is a 1-byte domain separator and every variable-width field
inside $\mathsf{ctx}$ is length-prefixed (e.g.\ $\mathrm{be8}(|A|) \| A$
for AAD $A$). Domain $\texttt{0x0A}$ is the single exception: it does not
take a nonce as input, since it is the derivation that \emph{produces} the
nonce (Definition~\ref{def:synthnonce}); its input shape is instead
$\texttt{0x0A}\|\mathrm{be8}(|A|)\|A\|M$.

\begin{definition}[Tag input]\label{def:taginput}
The authentication tag is always computed over the single quantity
\[
  \mathrm{TagIn}(A, \widetilde{B}) \;=\; \mathrm{be8}(|A|) \,\|\, A \,\|\, \widetilde{B},
\]
so that the complete domain-$\texttt{0x03}$ HMAC input is
$\texttt{0x03}\|N\|\mathrm{TagIn}(A,\widetilde{B})$. Every occurrence of the
tag in this paper — Table~\ref{tab:domains}, the wire format
(Section~\ref{sec:wire-format}), both $\mathsf{Enc}$ and $\mathsf{Dec}$
(Definition~\ref{def:aead-triple}), and the integrity proof
(Theorem~\ref{thm:int-ctxt}) — denotes this one quantity, so encryptor and
decryptor cannot diverge in the width of the length prefix.
\end{definition}

\begin{table}[h]
\centering
\begin{tabular}{@{}clll@{}}
\toprule
Domain & Function & Context & Output \\
\midrule
\texttt{0x00} & Noise position & $\mathrm{be5}(ct\_pos)$ & Boolean \\
\texttt{0x01} & Real-token addend & $\mathrm{be5}(real\_idx)$ & $[1, k_i - 1]$ \\
\texttt{0x02} & Noise probability & (none) & $[0.75, 0.99]$ \\
\texttt{0x03} & Auth tag & $\mathrm{TagIn}(aad, \widetilde{B})$ (Def.~\ref{def:taginput}) & 32 bytes \\
\texttt{0x04} & Noise character & $\mathrm{be5}(ct\_pos)$ & $[32, 127]$ \\
\texttt{0x05} & Noise-token addend & $\mathrm{be5}(ct\_pos)$ & $[1, k_i - 1]$ \\
\texttt{0x06} & Padding codepoint & $\mathrm{be4}(pad\_idx)$ & $[32, 126]$ \\
\texttt{0x07} & Keystream block & $\mathrm{be4}(block\_idx)$ & 32-byte block \\
\texttt{0x0A} & Synthetic nonce & $\mathrm{be8}(|A|) \| A \| M$ & 16 bytes \\
\bottomrule
\end{tabular}
\caption{Domain-separated HMAC derivation schedule. Every domain is keyed
by $\mathsf{sk}$ only. The nonce $N$ is applied uniformly by
$\mathrm{Derive}_d$ for $d\neq\texttt{0x0A}$ and is therefore omitted from
the Context column above.}
\label{tab:domains}
\end{table}

\begin{lemma}[Domain separation]\label{lem:domsep}
Write $\mathsf{enc}(d,N,\mathsf{ctx})$ for the HMAC-SHA256 input byte string
that $\mathrm{Derive}_d(\mathsf{sk},N,\mathsf{ctx})$ keys $\mathsf{sk}$
over, laid out as in Table~\ref{tab:domains}. Then $\mathsf{enc}$ is
injective on its \emph{entire} argument domain: for any two argument
tuples $(d_1,N_1,\mathsf{ctx}_1)\neq(d_2,N_2,\mathsf{ctx}_2)$ with each
field drawn from its declared range, the two resulting HMAC-SHA256 inputs
are distinct. The quantification is over all well-formed argument tuples,
\emph{not} merely over those this construction would itself generate; in
particular it covers tuples chosen adversarially, as
Theorem~\ref{thm:int-ctxt} requires. The statement is unconditional — it is
a property of the encoding alone, independent of $\mathsf{sk}$ and of any
computational assumption.
\end{lemma}

\begin{proof}
Two inputs with different domain bytes $d_1\neq d_2$ differ at byte
position $0$, regardless of any other field, so cross-domain distinctness
is unconditional. Fix a domain $d\neq\texttt{0x0A}$: bytes $1..16$ carry
the nonce, so if $N_1\neq N_2$ the two inputs already differ there; assume
therefore $N_1=N_2$, which reduces distinctness to injectivity of the
context encoding. For domains
$\{\texttt{0x00},\texttt{0x01},\texttt{0x02},\texttt{0x04},\texttt{0x05},\texttt{0x06},\texttt{0x07}\}$
the context is a fixed-width positional counter (or empty), trivially
injective;
for domain $\texttt{0x03}$, $\mathsf{ctx}=\mathrm{TagIn}(A,\widetilde{B})=\mathrm{be8}(|A|)\|A\|\widetilde{B}$,
and because $\mathrm{be8}(|A|)$ fixes $|A|$, the split between $A$ and
$\widetilde{B}$ is unambiguous, so distinct $(A,\widetilde{B})$ pairs never
collide. Within domain $\texttt{0x0A}$, whose input is
$\texttt{0x0A}\|\mathrm{be8}(|A|)\|A\|M$ with no nonce field at all, the
same length-prefix argument shows $\mathrm{be8}(|A|)\|A\|M$ is injective in
$(A,M)$. Every case is a statement about byte layout that holds for
arbitrary field values, so $\mathsf{enc}$ is injective on the whole domain:
any two distinct argument tuples, whatever their provenance, map to
distinct $\mathrm{HMAC}(\mathsf{sk},\cdot)$ inputs.
\end{proof}

\subsection{Noise Probability}
\label{sec:noiseprob}

The noise probability is realised entirely with fixed-width integer
arithmetic; no floating-point representation is used anywhere in this
construction. Fix the public integer constants
\[
  \theta_{\min} = \lfloor 0.75 \times 2^{64} \rfloor, \qquad
  \theta_{\max} = \lfloor 0.99 \times 2^{64} \rfloor,
\]
and let $\tau = \mathrm{Derive}_{0x02}(\mathsf{sk}, N, \varepsilon)[0:8]_{\mathrm{be64}} \in [0, 2^{64})$
be the raw 64-bit unsigned integer HMAC output (no division involved).
The per-message noise threshold is the 64-bit integer
\[
  \theta(N) = \theta_{\min} + \left\lfloor \frac{\tau \cdot (\theta_{\max} - \theta_{\min})}{2^{64}} \right\rfloor,
\]
computed with a single 128-bit intermediate product and one integer
division, so that $\theta(N)$ is bit-identical across all conforming
implementations regardless of platform or language. Every informal
reference elsewhere in this paper to a ``noise probability $p \in
[0.75, 0.99]$'' (e.g.\ Table~\ref{tab:domains}) denotes this integer
threshold, i.e.\ $p = \theta(N)/2^{64}$.

\subsection{Plaintext Padding}
\label{sec:padding}

Let $n = \lvert P \rvert$ (plaintext length in Unicode codepoints). Define
\[
  B = \max(16,\ 2^{\lceil \log_2(n+1) \rceil}), \qquad P_d = [n \gg 8,\ n \wedge \mathrm{0xFF}] \| P \| \mathbf{pad},
\]
where $\mathbf{pad} = (\mathbf{pad}_0, \ldots, \mathbf{pad}_{B-n-1})$ is a
$(B-n)$-codepoint sequence with
\[
  \mathbf{pad}_{i} = \bigl(\mathrm{Derive}_{0x06}(\mathsf{sk}, N, \mathrm{be4}(i))[0{:}4]_{\mathrm{be32}} \bmod 95\bigr) + 32 \ \in [32,126].
\]
The total padded length is $R = n + 2 + (B-n) = B+2$ codepoints; we call
$B$ the \emph{padding bucket} of the message and $R$ the \emph{real-token
count}. Because $B$ depends only on $n$ (never on the codepoint values
themselves), two messages of equal codepoint length always share the same
padding bucket and the same real-token count.

\subsection{Padding Profiles}
\label{sec:pad-profiles}

The map $n\mapsto B$ above is the \emph{bucket} profile, and it is the
normative default. Because it is the sole source of the length-hiding
property proved in Section~\ref{sec:length-hiding}
(Remark~\ref{rem:coarsening-not-inflation}), we specify it as a replaceable
parameter rather than as a hard-wired constant. A \emph{padding profile} is
any function $B(\cdot)$ from codepoint counts to block sizes satisfying
$B(n)>n$ throughout the message space. This paper specifies three:
\begin{description}
  \item[$\mathsf{bucket}$ (default).]
        $B(n)=\max\bigl(16,\ 2^{\lceil\log_2(n+1)\rceil}\bigr)$, with $13$
        reachable values $\{2^4,2^5,\ldots,2^{16}\}$.
  \item[$\mathsf{coarse}(g)$.] $B(n)=2^{\,4+g\lceil(e(n)-4)/g\rceil}$ where
        $e(n)=\max\bigl(4,\lceil\log_2(n+1)\rceil\bigr)$ and $g$ divides
        $12$: the same ladder thinned by a stride $g$, with $12/g+1$
        reachable values. $\mathsf{coarse}(1)=\mathsf{bucket}$, and
        $\mathsf{coarse}(12)$ leaves only $\{2^4,2^{16}\}$.
  \item[$\mathsf{frame}(F)$.] $B(n)=F$ for a fixed
        $F\in\{2^4,\ldots,2^{16}\}$, on the restricted message space $n<F$:
        every message pads to one common block size, so exactly one value is
        reachable.
\end{description}

All three profiles take values in the same $13$-element set, so the set of
legal real-token counts $R=B+2$ — and hence the set of legal ciphertext
lengths — is identical for all of them. A decryptor consequently needs
\emph{no} knowledge of the sender's profile: it recovers $R$ from the token
count (Section~\ref{sec:wire-format}) and the true length $n$ from the
2-codepoint length prefix, so any conforming decryptor accepts any
profile's output unchanged, and any well-formedness check on $R$ remains a
check against the same fixed set. The profile is a deployment parameter
agreed out of band, exactly like $\mathbf{k}$ and $K$; it is never
transmitted and does not appear in the wire format.

Corollary~\ref{cor:length-leak} is stated over an arbitrary profile: length
leakage is at most $\log_2$(number of reachable values) bits — about $3.70$
bits under $\mathsf{bucket}$, and \emph{exactly zero} under
$\mathsf{frame}(F)$, at the cost of every message paying the $F$-sized
ciphertext. Deployments whose threat model is dominated by traffic analysis
should select $\mathsf{frame}(F)$ for the largest $F$ their bandwidth budget
tolerates; this, and not the expansion factor, is the knob that trades
bandwidth for length confidentiality
(Remark~\ref{rem:coarsening-not-inflation}).

\subsection{Token Construction}
\label{sec:token-construction}

For each padded codepoint $c$, define the following primitives, each an
explicit reduction of a 32-byte $\mathrm{Derive}_d$ output into its target
range:

\begin{definition}[Token-construction primitives]\label{def:token-prims}
For a key element $k$, position counters $ct\_pos$, $real\_idx$:
\begin{align*}
  \texttt{is\_noise}(ct\_pos) &= \mathbf{1}\bigl[\,\mathrm{Derive}_{0x00}(\mathsf{sk}, N, \mathrm{be5}(ct\_pos))[0{:}8]_{\mathrm{be64}} < \theta(N)\,\bigr], \\
  \texttt{noise\_char}(ct\_pos) &= \bigl(\mathrm{Derive}_{0x04}(\mathsf{sk}, N, \mathrm{be5}(ct\_pos))[0{:}4]_{\mathrm{be32}} \bmod 96\bigr) + 32, \\
  \texttt{noise\_addend}(ct\_pos, k) &= \bigl(\mathrm{Derive}_{0x05}(\mathsf{sk}, N, \mathrm{be5}(ct\_pos))[0{:}4]_{\mathrm{be32}} \bmod (k-1)\bigr) + 1, \\
  \texttt{real\_addend}(real\_idx, k) &= \bigl(\mathrm{Derive}_{0x01}(\mathsf{sk}, N, \mathrm{be5}(real\_idx))[0{:}4]_{\mathrm{be32}} \bmod (k-1)\bigr) + 1.
\end{align*}
\end{definition}

\begin{remark}[Modulo bias]\label{rem:modbias}
Each finite-range reduction above (and the padding-codepoint reduction of
Section~\ref{sec:padding}) draws a value pseudo-uniformly from
$[0,2^{32})$ and reduces it modulo a small integer $m$. Writing
$2^{32}=am+r$ with $0\le r<m$, exactly $r$ residues occur with probability
$(a+1)/2^{32}$ and the remaining $m-r$ with probability $a/2^{32}$, so the
statistical distance from uniform on $\mathbb{Z}_m$ is at most $m/2^{32}$:
below $2^{-25}$ for the padding codepoint ($m=95$) and the noise character
($m=96$), and below $2^{-8}$ for the addends ($m=k-1<2^{24}$).

We do \emph{not} claim that this bias is ``dominated by'' the PRF advantage
of HMAC-SHA256; that comparison is not meaningful, since a small PRF
advantage says nothing about the uniformity of a reduced output. The
correct statement is that no result in Section~\ref{sec:security} requires
any of these reduced values to be uniform. Confidentiality rests entirely
on the one-time keystream $\mathrm{ks}(N)$ that masks the serialised token
blob (Section~\ref{sec:wire-format}), and integrity entirely on the tag;
the addends, noise characters and padding codepoints are computed
\emph{underneath} that mask and are never observed by the adversary in the
clear. The one correctness property that does depend on them,
$\gcd(a,k)=1$, follows from the range $a\in[1,k-1]$ and holds for every
draw regardless of bias. No rejection sampling is therefore required — not
because the bias is negligible relative to some other quantity, but because
nothing is claimed that it could invalidate.

The noise threshold $\theta(N)$ (Section~\ref{sec:noiseprob}) deserves an
explicit correction here: it is \emph{not} bias-free. Although it uses exact
integer arithmetic and no modulo reduction, the fixed-point scaling
$\theta_{\min}+\lfloor\tau(\theta_{\max}-\theta_{\min})/2^{64}\rfloor$ maps
$2^{64}$ equally likely draws of $\tau$ onto approximately
$0.24\cdot2^{64}$ outputs, i.e.\ about $4.17$ preimages per output. Since
each output receives either $4$ or $5$ preimages, individual outputs
deviate from uniform on the range by up to $+20\%$ in relative terms, and
the total variation distance from the uniform distribution on
$[\theta_{\min},\theta_{\max}]$ is $\approx0.033$. This is harmless — the
only property required of $\theta(N)$ is that it land in the public
interval $[\theta_{\min},\theta_{\max}]$, which it does by construction,
and that $\mathsf{Dec}$ recompute it identically, which it does — but it
should not be described as exempt. The prime-index sampling of
Remark~\ref{rem:sampling} is genuinely bias-free, by rejection.

Finally, the boundary of this argument is narrow and worth stating, because
it is a property of the composition rather than of the reductions
themselves. Any variant that exposes a reduced value before it is masked,
that drops the keystream mask, or that reuses a nonce across two messages
makes these biases immediately relevant, and the analysis above would then
have to be redone with explicit bias terms. Implementers must not treat
``no rejection sampling needed'' as a property of the derivations in
isolation.
\end{remark}

Real and noise tokens are integers $c\cdot k + a$ for codepoint $c$
(or a noise character), key element $k=\mathbf{k}[i \bmod K]$, and addend
$a$ from Definition~\ref{def:token-prims}. Because $k$ is prime and every
addend lies in $[1,k-1]$, $\gcd(a,k)=1$ and hence $\gcd(\text{token},k)=1$
for every token, real or noise.

To keep ciphertext length independent of the message-derived nonce's
particular noise realisation (Section~\ref{sec:wire-format}), the number
of \emph{consecutive} noise tokens emitted before each real token is capped
at a fixed constant $\mathrm{MAX\_NOISE\_RUN}=19$, and the total token
count is padded, with additional filler tokens structurally identical to
genuine noise tokens, up to a fixed ceiling of exactly
$R\cdot(\mathrm{MAX\_NOISE\_RUN}+1)$ tokens per message:

\begin{lstlisting}[language=Python, caption={Token emission (encoding)}]
for c in padded:                     # R codepoints total
    noise_run = 0
    while noise_run < MAX_NOISE_RUN and is_noise(ct_pos):
        k = key[real_idx % K]
        emit noise_char(ct_pos) * k + noise_addend(ct_pos, k)
        ct_pos += 1; noise_run += 1
    k = key[real_idx % K]
    emit c * k + real_addend(real_idx, k)
    ct_pos += 1; real_idx += 1

# Ceiling padding: append filler tokens (identical construction to
# genuine noise tokens) until exactly R * (MAX_NOISE_RUN + 1) tokens
# have been emitted in total.
ceiling = R * (MAX_NOISE_RUN + 1)
while len(tokens) < ceiling:
    k = key[real_idx % K]
    emit noise_char(ct_pos) * k + noise_addend(ct_pos, k)
    ct_pos += 1
\end{lstlisting}

Because the ceiling depends only on $R$ (hence only on the padding bucket
$B$), the total token count $N_{\mathrm{tok}} = R\cdot(\mathrm{MAX\_NOISE\_RUN}+1)$
is, by construction, a deterministic function of the padding bucket
alone — regardless of $\mathsf{sk}$, $N$, $A$, or the specific message
(Lemma~\ref{lem:length-det}).

\subsection{Synthetic Nonce}
\label{sec:synthnonce}

\begin{definition}[Synthetic nonce]\label{def:synthnonce}
For $\mathsf{sk}\in\{0,1\}^{256}$, $A\in\{0,1\}^{*}$, $M\in\mathcal{M}$
(encoded as UTF-8 bytes),
\[
  N(\mathsf{sk}, A, M) \;=\; \mathrm{HMAC}\bigl(\mathsf{sk},\;
    \texttt{0x0A} \,\|\, \mathrm{be8}(|A|) \,\|\, A \,\|\, M\bigr)[0{:}16].
\]
\end{definition}

This is a synthetic IV in the manner of RFC~5297 SIV / AES-GCM-SIV: it is
a deterministic function of the key and the input, not sampled
independently at random, so re-encrypting an identical $(A,M)$ pair always
reproduces the identical nonce (and hence the identical ciphertext) —
the standard, disclosed message-equality trade-off of deterministic,
misuse-resistant encryption (Definition~\ref{def:ind-cpa-det}).

\subsection{Wire Format}
\label{sec:wire-format}

Each token is serialised in a fixed-width 8-byte field,
\[
  B = \mathrm{be8}(\mathrm{token}[0]) \| \mathrm{be8}(\mathrm{token}[1]) \| \cdots \| \mathrm{be8}(\mathrm{token}[N_{\mathrm{tok}}-1]),
\]
so $|B| = 8N_{\mathrm{tok}}$ bytes, with $N_{\mathrm{tok}} = R\cdot(\mathrm{MAX\_NOISE\_RUN}+1)$
as fixed in Section~\ref{sec:token-construction}. $B$ is XOR-masked with a
keystream derived from domain $\texttt{0x07}$:
\[
  \widetilde{B} \;=\; B \,\oplus\, \mathrm{ks}(N)\bigl[0 : |B|\bigr],
  \qquad
  \mathrm{ks}(N) \;=\; \mathrm{Derive}_{0x07}(\mathsf{sk}, N, \mathrm{be4}(0)) \,\|\,
  \mathrm{Derive}_{0x07}(\mathsf{sk}, N, \mathrm{be4}(1)) \,\|\, \cdots,
\]
generating only as many 32-byte blocks as needed to cover $|B|$ bytes
(truncating the final block). The authentication tag is
\[
  T = \mathrm{Derive}_{0x03}\bigl(\mathsf{sk},\ N,\ \mathrm{TagIn}(A,\widetilde{B})\bigr),
\]
a 32-byte HMAC-SHA256 value binding the (possibly empty) associated data
$A$. The ciphertext is
\[
  C = N \| \widetilde{B} \| T,
\]
a 16-byte nonce, the masked token blob, and the 32-byte tag. A decryptor
rejects any input with $|C| < 48$ bytes outright (parse-time floor); the
true minimum length of a well-formed ciphertext is larger, since every
message — including the empty string — pads to at least $B=16$ codepoints
($R=18$ real tokens), giving
$N_{\mathrm{tok}} = 18\cdot20 = 360$, $|B|=2880$ bytes, and
$|C| = 16+2880+32 = 2928$ bytes at minimum. In general
$|C| = 48 + 160\,(B(|M|)+2)$ bytes, a closed form in the padding bucket
alone.

\begin{lemma}[Ciphertext length is a deterministic function of the padding bucket]
\label{lem:length-det}
For any $\mathbf{k}$, $\mathsf{sk}$, $A$, $M$,
\[
  |C| \;=\; 16 \;+\; 8\,R\,(\mathrm{MAX\_NOISE\_RUN}+1) \;+\; 32,
  \qquad R = B(|M|) + 2,
\]
where $B(\cdot)$ is the padding-bucket function of
Section~\ref{sec:padding}. In particular $|C|$ depends on $M$ only through
$B(|M|)$ — never through $\mathsf{sk}$, $N$, $A$, or the specific
codepoints of $M$.
\end{lemma}

\begin{proof}
Immediate from the construction of Section~\ref{sec:token-construction}:
$N_{\mathrm{tok}}$ is fixed, by the ceiling-padding step, to exactly
$R\cdot(\mathrm{MAX\_NOISE\_RUN}+1)$ tokens regardless of the natural
noise realisation, and $R = B(|M|)+2$ depends only on the codepoint
\emph{count} of $M$ (Section~\ref{sec:padding}). The wire format then adds
a constant 16-byte nonce and 32-byte tag, and masking (XOR) does not
change $|B|$.
\end{proof}

\subsection{The NAPQES Algorithm Triple}
\label{sec:algorithm-triple}

\begin{definition}[NAPQES AEAD algorithm triple]\label{def:aead-triple}
NAPQES is the triple $\Pi = (\mathsf{KeyGen}, \mathsf{Enc}, \mathsf{Dec})$ over
the following typed spaces. The symbol $\Pi$ denotes this scheme, and only
this scheme, throughout Section~\ref{sec:security}.
\begin{itemize}
  \item Key space $\mathcal{K}\times\{0,1\}^{256}$: $\mathcal{K}$ is the set
        of ordered $K$-tuples of distinct elements of $\mathcal{P}$
        (Section~\ref{sec:keygen}).
  \item AAD space $\mathcal{A} = \{0,1\}^{*}$: the AAD length is encoded
        as $\mathrm{be8}(|A|)$ throughout (Section~\ref{sec:domains}), so
        any $|A|<2^{64}$ bytes is supported unambiguously, with no
        practical ceiling.
  \item Message space $\mathcal{M}$: finite sequences of Unicode
        \emph{scalar values} — that is, codepoints in
        $[\mathtt{U+0000},\mathtt{U+10FFFF}]\setminus[\mathtt{U+D800},\mathtt{U+DFFF}]$
        — of length at most
        $\mathrm{MAX\_PLAINTEXT\_CODEPOINTS}=2^{16}-1$, the
        limit imposed by the 2-codepoint big-endian length prefix written
        during padding (Section~\ref{sec:padding}). The surrogate range is
        excluded because the synthetic-nonce derivation
        (Definition~\ref{def:synthnonce}) hashes the UTF-8 encoding of $M$,
        and UTF-8 is undefined on unpaired surrogates; $\mathcal{M}$ is thus
        exactly the set of strings that admit a UTF-8 encoding. Nothing
        weaker is needed and nothing stronger is assumed: the arithmetic
        layer itself operates on the integer codepoint values and is
        indifferent to which subset of $[0,2^{21})$ they range over.
  \item Ciphertext space $\mathcal{C}=\{0,1\}^{*}$: every output of
        $\mathsf{Enc}$ has length $16+8R(\mathrm{MAX\_NOISE\_RUN}+1)+32$
        bytes for $R=B(|M|)+2$ (Lemma~\ref{lem:length-det}), that is, in
        closed form,
        \[
          |C| \;=\; 48 + 160\bigl(B(|M|)+2\bigr)\ \text{bytes};
        \]
        the parse-time
        rejection floor is 48 bytes, but the true minimum length of a
        well-formed ciphertext is $48+160\cdot18=2928$ bytes.
\end{itemize}

\begin{remark}[The token field is wide enough for the whole of $\mathcal{M}$]
\label{rem:token-fit}
Every emitted token has the form $c\cdot k_j+a$ with $c$ a padded codepoint,
$k_j\in\mathcal{P}$ and $a\in[1,k_j-1]$
(Section~\ref{sec:token-construction}), so
$c\cdot k_j+a\le(c+1)k_j-1$. Taking the extreme values admitted by
$\mathcal{M}$ and $\mathcal{P}$ — $c=\mathtt{0x10FFFF}$ and
$k_j=9{,}899{,}999$, the largest prime in $\mathcal{P}$ — gives a maximum
token of $11{,}029{,}707{,}685{,}887\approx2^{43.33}$. The fixed 8-byte
big-endian token field therefore never overflows, with a margin of a factor
of over $1.6\times10^6$. This bound is unconditional: it holds for every
$M\in\mathcal{M}$, every $\mathbf{k}$, and every noise realisation, so
Lemma~\ref{lem:length-det} and the serialisation are well-defined on the
entire message space.
\end{remark}

\begin{remark}[Implementation domains]\label{rem:impl-domain}
The Python and Rust reference implementations realise $\mathcal{M}$ as
specified: Rust's \texttt{\&str} excludes surrogates by construction, and
the Python port rejects them explicitly. The C port
(\texttt{C/napqes.c}) takes a NUL-terminated \texttt{char\,*} and maps each
\emph{byte} to one codepoint, so it realises $\mathcal{M}$ faithfully only
for ASCII input; on non-ASCII input it produces a different, self-consistent
ciphertext that the other two ports will not reproduce. This is an
implementation limitation of that port, not a property of the scheme, and
it is recorded in Section~\ref{sec:caveats}.
\end{remark}

\paragraph{$\mathsf{KeyGen}(1^{\lambda}) \to (\mathbf{k},\mathsf{sk})$.}
As in Definition~\ref{def:keygen}.

\paragraph{$\mathsf{Enc}(\mathbf{k}, \mathsf{sk}, A, M) \to C$.}
(1) compute $N=N(\mathsf{sk},A,M)$ (Definition~\ref{def:synthnonce}); (2)
pad $M$ to $P_d$ (Section~\ref{sec:padding}); (3) run the token-emission
loop of Section~\ref{sec:token-construction} over $P_d$ using
$(\mathsf{sk},N)$ to obtain the token vector
$(\mathrm{token}[0],\ldots,\mathrm{token}[N_{\mathrm{tok}}-1])$; (4)
serialise $B$; (5) mask $\widetilde{B}=B\oplus\mathrm{ks}(N)$; (6) compute
$T=\mathrm{Derive}_{0x03}(\mathsf{sk},N,\mathrm{TagIn}(A,\widetilde{B}))$
(Definition~\ref{def:taginput});
(7) output $C=N\|\widetilde{B}\|T$. The public encryption oracle is
$\mathsf{Encrypt}(\mathbf{k},\mathsf{sk},A,M)\triangleq\mathsf{Enc}(\mathbf{k},\mathsf{sk},A,M)$
— a \emph{pure function} of $(A,M)$, since $N$ is computed synthetically
rather than sampled.

\paragraph{$\mathsf{Dec}(\mathbf{k}, \mathsf{sk}, A, C) \to M \text{ or } \bot$.}
(1) if $|C|<48$, output $\bot$; (2) parse $C=N\|\widetilde{B}\|T'$ with $N$
the leading 16 bytes and $T'$ the trailing 32 bytes; (3) let
$N_{\mathrm{tok}}=(|C|-48)/8$ and, if this is not a non-negative integer or
not divisible by $\mathrm{MAX\_NOISE\_RUN}+1$, output $\bot$; otherwise set
$R=N_{\mathrm{tok}}/(\mathrm{MAX\_NOISE\_RUN}+1)$; (4) compute
$T=\mathrm{Derive}_{0x03}(\mathsf{sk},N,\mathrm{TagIn}(A,\widetilde{B}))$
— the same quantity, with the same $\mathrm{be8}$ length prefix, that
$\mathsf{Enc}$ step (6) computes (Definition~\ref{def:taginput});
(5) if $T\neq T'$ (compared in constant time), output $\bot$;
(6) if $R\notin\{B+2 : B\in\{2^4,\ldots,2^{16}\}\}$, output $\bot$
(Remark~\ref{rem:dec-structural});
(7) else unmask $B=\widetilde{B}\oplus\mathrm{ks}(N)$, deserialise the
$N_{\mathrm{tok}}$ fixed-width tokens, and invert the token-emission loop
using the same $(\mathsf{sk},N)$-derived noise-position oracle: skip at most
$\mathrm{MAX\_NOISE\_RUN}$ consecutive noise positions, then read one real
token $t$ and recover the padded codepoint as
$(t-a)/k_{(\mathrm{real\_idx}\bmod K)+1}$, where $a$ is the
$\mathrm{real\_idx}$-th addend
$\mathrm{Derive}_{0x01}(\mathsf{sk},N,\mathrm{real\_idx})$ and
$\mathrm{real\_idx}$ counts real tokens from $0$ — exactly the prime and
addend $\mathsf{Enc}$ used at that index; stop once $R$ real tokens have
been extracted, discarding any trailing filler; (8) read the 2-codepoint
big-endian length prefix $n$ and, if $n>R-2$, output $\bot$
(Remark~\ref{rem:dec-structural}); (9) output the $n$ codepoints following
the prefix.
\end{definition}

\begin{remark}[The post-authentication structural checks]
\label{rem:dec-structural}
Steps (6) and (8) are performed \emph{after} the tag comparison of step (5),
so they are never reachable on an input an adversary without $\mathsf{sk}$
could construct, and they play no part in the INT-CTXT argument. They are
specified nonetheless because step (7) is not total without them: $R$ is
recovered from $|C|$ and $n$ from the decrypted buffer, so a malformed but
correctly tagged ciphertext — which a key holder can trivially produce —
would otherwise drive an out-of-range read or an allocation sized by a field
the caller did not choose. Divisibility by $\mathrm{MAX\_NOISE\_RUN}+1$
alone does not imply step (6): $R=20$ passes step (3) and fails step (6).
All three reference implementations enforce both checks, and vectors
\texttt{W-N06}--\texttt{W-N08} of Section~\ref{sec:kats} pin them.
\end{remark}

\begin{proposition}[Correctness]\label{prop:correctness}
For every $\lambda$, every $(\mathbf{k},\mathsf{sk})\gets\mathsf{KeyGen}(1^\lambda)$,
every $A\in\mathcal{A}$, $M\in\mathcal{M}$:
\[
  \mathsf{Dec}\bigl(\mathbf{k},\mathsf{sk},\,A,\,\mathsf{Enc}(\mathbf{k},\mathsf{sk},A,M)\bigr) = M.
\]
\end{proposition}

\begin{proof}
Let $C=N\|\widetilde{B}\|T$ be the output of
$\mathsf{Enc}(\mathbf{k},\mathsf{sk},A,M)$. We first check that the two tag
inputs coincide, rather than assuming tag agreement.

By Lemma~\ref{lem:length-det}, $|C|=16+8R(\mathrm{MAX\_NOISE\_RUN}+1)+32$
with $R=B(|M|)+2\ge1$, so $|C|\ge2928>48$ and
$N_{\mathrm{tok}}=(|C|-48)/8=R(\mathrm{MAX\_NOISE\_RUN}+1)$ is a
non-negative integer divisible by $\mathrm{MAX\_NOISE\_RUN}+1$: steps (1)
and (3) of $\mathsf{Dec}$ do not reject, and step (3) recovers exactly the
$R$ used by $\mathsf{Enc}$. Step (2) parses the leading 16 bytes as $N$ and
the trailing 32 as $T'=T$, leaving the middle bytes as exactly the
$\widetilde{B}$ that $\mathsf{Enc}$ emitted; in particular $N$ is
\emph{read from the ciphertext}, not re-derived, so no circularity with the
synthetic-nonce derivation arises. Hence $\mathsf{Dec}$ step (4) evaluates
$\mathrm{Derive}_{0x03}$ on $\mathrm{TagIn}(A,\widetilde{B})$ with the same
$\mathsf{sk}$, the same $N$, the same caller-supplied $A$, and the same
$\widetilde{B}$ as $\mathsf{Enc}$ step (6) did
(Definition~\ref{def:taginput}). The two tag inputs are byte-identical, so
$T=T'$ and step (5) does not reject. Since $R=B(|M|)+2$ with
$B(|M|)\in\{2^4,\ldots,2^{16}\}$, step (6) does not reject either, and since
the length prefix written by $\mathsf{Enc}$ satisfies $n=|M|\le B(|M|)=R-2$,
neither does step (8).

Step (7) then unmasks with $\mathrm{ks}(N)$, the same keystream
$\mathsf{Enc}$ masked with (a function of $\mathsf{sk}$ and the recovered
$N$ alone), recovering $B$ and hence the token vector. $\mathsf{Enc}$'s
token-emission loop and $\mathsf{Dec}$'s loop-inversion step both compute
the noise-position oracle, addends, and padding codepoints as pure,
deterministic functions of $(\mathsf{sk},N,\mathrm{position})$, so with the
same $\mathsf{sk}$ and $N$ they compute identical intermediate values at
every position; emission and inversion are therefore exact inverses, and
stripping the length prefix and padding yields $M$.
\end{proof}

\begin{remark}[Valid ciphertexts are not unique per $(A,M)$]
\label{rem:not-unique}
Because $\mathsf{Dec}$ reads $N$ from the ciphertext instead of re-deriving
it from $(A,M)$, the set of $C$ that decrypt to a given $M$ under a given
$A$ is larger than $\{\mathsf{Enc}(\mathbf{k},\mathsf{sk},A,M)\}$: anyone
holding $\mathsf{sk}$ may pick an arbitrary $N'\neq N$, rebuild
$\widetilde{B}'$ and $T'$ under it, and obtain a second, equally valid
ciphertext for the same plaintext. $\mathsf{Enc}$ is a function, but
$\mathsf{Dec}$ is not injective on valid ciphertexts, so NAPQES is
\emph{not} tidy in the sense of Namprempre--Rogaway--Shrimpton, and
ciphertext equality must not be used as a proxy for plaintext equality by
any protocol built on it.

This costs nothing in the security analysis. INT-CTXT
(Definition~\ref{def:int-ctxt}) measures freshness over the submitted pair
$(c,A)$, not over the underlying plaintext, so an alternative encoding of an
already-queried message is a legitimate forgery attempt and is counted as
one; Theorem~\ref{thm:int-ctxt} bounds it. Re-deriving $N$ inside
$\mathsf{Dec}$ and rejecting on mismatch would restore uniqueness, but it
would also make $\mathsf{Dec}$ depend on the plaintext it is trying to
recover, and it is not done.
\end{remark}

\begin{remark}[Authentication-failure output $\bot$]
The formal $\bot$ output of $\mathsf{Dec}$ is realised in all three
reference implementations as a raised exception (\texttt{ValueError} in
Python/Rust, an equivalent error code in C) rather than a sentinel return
value — a standard, semantically equivalent realisation of $\bot$ (cf.\
most software AEAD APIs, e.g.\ \texttt{cryptography.hazmat}).
\end{remark}

% ─────────────────────────────────────────────────────────────────────────────
\section{Security Analysis}
\label{sec:security}
% ─────────────────────────────────────────────────────────────────────────────

Every theorem in this section is unconditional under the standard,
uniform-key PRF assumption on HMAC-SHA256 stated below: because
$\mathsf{sk}$ is sampled uniformly and independently of $\mathbf{k}$, every
reduction constructed in this section may sample $\mathbf{k}$ on its own
(to run the arithmetic layer) while forwarding every HMAC call to an
external oracle keyed by the challenger's $\mathsf{sk}$ — the reduction
never needs to extract $\mathsf{sk}$ from $\mathbf{k}$, or vice versa, so
no simulation gap and no non-standard key-distribution assumption arises.
NAPQES does not define a standalone MAC verification key separate from
$\mathsf{sk}$, so EUF-CMA unforgeability of the tag and INT-CTXT
unforgeability of the ciphertext coincide here.

\subsection{PRF Advantage and Truncation}

\begin{definition}[HMAC-SHA256 PRF advantage]\label{def:prf-adv}
Let $F=\mathrm{HMAC\text{-}SHA256}:\{0,1\}^{*}\times\{0,1\}^{*}\to\{0,1\}^{256}$.
For a PPT distinguisher $\mathcal{B}$ with oracle access to either
$F(\kappa,\cdot)$ for $\kappa\xleftarrow{\$}\{0,1\}^{256}$, or a function
$\rho:\{0,1\}^{*}\to\{0,1\}^{256}$ sampled uniformly from all functions of
that type,
\[
  \mathrm{Adv}^{\mathrm{PRF}}_{F}(\mathcal{B})
  \;=\;
  \Bigl|\Pr\bigl[\mathcal{B}^{F(\kappa,\cdot)}=1\bigr]-\Pr\bigl[\mathcal{B}^{\rho(\cdot)}=1\bigr]\Bigr|.
\]
\end{definition}

\begin{assumption}[Uniform-key HMAC-SHA256 PRF]\label{assum:hmac-prf}
$\mathrm{Adv}^{\mathrm{PRF}}_{\mathrm{HMAC\text{-}SHA256}}(\mathcal{B})$
is negligible in the security parameter for every PPT $\mathcal{B}$,
for $\kappa$ drawn uniformly from $\{0,1\}^{256}$ (matching how
$\mathsf{sk}$ is sampled, Definition~\ref{def:keygen}).
\end{assumption}

\begin{lemma}[Truncation preserves the PRF property]\label{lem:trunc-prf}
Let $F:\{0,1\}^{256}\times\{0,1\}^*\to\{0,1\}^{256}$ be a keyed function
with 32-byte output and fix a truncation length $\ell$ in \emph{bytes},
$1\le\ell\le32$; write $\mathrm{trunc}_\ell(x)=x[0{:}\ell]$, so that
$\mathrm{trunc}_\ell\circ F$ has output length $8\ell$ bits. Then for every
adversary $\mathcal{B}$ against $\mathrm{trunc}_\ell\circ F$ there is an
adversary $\mathcal{B}'$ against $F$, running in essentially the same time
and making the same number of queries, with
\[
  \mathrm{Adv}^{\mathrm{PRF}}_{\mathrm{trunc}_\ell\circ F}(\mathcal{B})
  \;=\;
  \mathrm{Adv}^{\mathrm{PRF}}_{F}(\mathcal{B}').
\]
\end{lemma}

\begin{proof}
Define $\mathcal{B}'$ to be the adversary that runs $\mathcal{B}$, forwards
each of $\mathcal{B}$'s queries to its own oracle, truncates each 32-byte
answer to its first $\ell$ bytes before handing it back to $\mathcal{B}$,
and finally outputs whatever $\mathcal{B}$ outputs. $\mathcal{B}'$ makes
exactly as many queries as $\mathcal{B}$ and adds only the cost of the
truncations.

If $\mathcal{B}'$'s oracle is $F(\kappa,\cdot)$ then the view it presents to
$\mathcal{B}$ is exactly $(\mathrm{trunc}_\ell\circ F)(\kappa,\cdot)$. If
$\mathcal{B}'$'s oracle is a uniformly random $\rho$ with 32-byte outputs,
then $\mathrm{trunc}_\ell\circ\rho$ is a uniformly random function with
$\ell$-byte outputs — each answer is a prefix of an independent uniform
32-byte string, hence uniform on $\{0,1\}^{8\ell}$, and answers at distinct
inputs remain independent — which is exactly the ideal object in
$\mathcal{B}$'s game. Both probabilities therefore transfer verbatim, and
the two advantages are equal. Note the direction of the reduction: it is
$\mathcal{B}$ that is simulated by $\mathcal{B}'$, so the equality is
between the advantage of $\mathcal{B}$ in one game and that of the
\emph{derived} adversary $\mathcal{B}'$ in the other, not between the
advantages of a single adversary in two different games.
\end{proof}

This covers every truncated derivation used by this construction (nonce
$\ell=16$, addend and padding codepoint $\ell=4$, noise threshold
$\ell=8$) under Assumption~\ref{assum:hmac-prf} alone. The symbol $\ell$ is
local to this lemma; it is unrelated to the noise-threshold draw $\tau$
appearing in the definition of $\theta(N)$ in Section~\ref{sec:noiseprob}.

\subsection{IND-CPA-det Security}
\label{sec:ind-cpa}

Because $\mathsf{Encrypt}(\mathbf{k},\mathsf{sk},A,M)$ is a pure function
of $(A,M)$ (Section~\ref{sec:algorithm-triple}), the standard,
unrestricted left-or-right IND-CPA game is not achievable by any
deterministic scheme: an adversary who queries the encryption oracle on
one of its own two challenge messages trivially learns the challenge bit
by comparing the result to the challenge ciphertext. That this game is
unachievable is not a defect of the construction but a structural property
of determinism, shared by every deterministic scheme including the
synthetic-IV designs of Rogaway and Shrimpton~\cite{rogaway2006}. The
notion we prove in its place is the restricted-query game below.
Remark~\ref{rem:dae-scope} states how that notion relates to, and where it
falls short of, the deterministic authenticated-encryption (DAE) notion
of~\cite{rogaway2006}.

\begin{definition}[IND-CPA-det security]\label{def:ind-cpa-det}
For a PPT adversary $\mathcal{A}=(\mathcal{A}_1,\mathcal{A}_2)$, define
$\mathbf{Exp}^{\mathrm{ind\text{-}cpa\text{-}det}}_{\mathcal{A}}(\lambda)$:
\begin{enumerate}
  \item $(\mathbf{k},\mathsf{sk})\gets\mathsf{KeyGen}(1^\lambda)$.
  \item $\mathcal{A}_1$ is given adaptive oracle access to
        $\mathsf{Encrypt}(\mathbf{k},\mathsf{sk},\cdot,\cdot)$ and, after at
        most $q$ total queries (counted jointly with step 4), outputs
        $(A^*,m_0,m_1,\mathsf{st})$ with $|m_0|=|m_1|$ (codepoint count).
  \item The challenger samples $b\xleftarrow{\$}\{0,1\}$ and returns
        $c^*\gets\mathsf{Encrypt}(\mathbf{k},\mathsf{sk},A^*,m_b)$.
  \item $\mathcal{A}_2(\mathsf{st},c^*)$, with continued oracle access,
        outputs a bit $b'$.
  \item The experiment outputs $1$ iff $b'=b$.
\end{enumerate}
$\mathcal{A}$ must additionally respect two restrictions, or the experiment
aborts with $\mathcal{A}$ losing: (i) \emph{no repeated queries} — $\mathcal{A}$
never queries the encryption oracle twice on an identical $(A,M)$ pair, in
either phase; (ii) \emph{no challenge-message queries} — writing
$(A^*,m_0,m_1)$ for $\mathcal{A}_1$'s output, $\mathcal{A}$ never queries
the encryption oracle, in either phase, on $(A^*,m_0)$ or $(A^*,m_1)$.
Restriction (ii) is the one forced by determinism, as argued above.
Restriction (i) costs no generality: because $\mathsf{Encrypt}$ is a pure
function, an adversary that repeats a query already holds the answer, so
any $\mathcal{A}$ violating (i) is converted into one respecting it, with
identical view and identical advantage, by caching its own
query/response pairs and answering repeats locally.
The advantage of $\mathcal{A}$ is
$\mathrm{Adv}^{\mathrm{IND\text{-}CPA\text{-}det}}_{\Pi}(\mathcal{A}) =
|\Pr[\mathbf{Exp}^{\mathrm{ind\text{-}cpa\text{-}det}}_{\mathcal{A}}(\lambda)=1]-\tfrac12|$.
\end{definition}

\begin{remark}[Scope: this is not a DAE claim]\label{rem:dae-scope}
Definition~\ref{def:ind-cpa-det} is a \emph{single-challenge, left-or-right}
notion. The DAE notion of~\cite{rogaway2006}, which AES-GCM-SIV\cite{rfc8452}
and SIV~\cite{rfc5297} target, is stronger in two respects: it asks that
\emph{every} ciphertext returned over the whole interaction be
indistinguishable from uniformly random bits of the same length, not that
one designated challenge hide one bit; and it bundles that requirement with
ciphertext integrity into a single game. We prove neither of those here.
What Theorem~\ref{thm:ind-cpa} and Theorem~\ref{thm:int-ctxt} establish is
the pair of separate, weaker statements, composed into
Theorem~\ref{thm:ind-cca} in the usual way. Accordingly, references in this
paper to AES-GCM-SIV and to ``misuse resistance'' denote a shared design
philosophy — a nonce synthesised from the message so that nonce reuse is
not a caller-controlled failure mode — and must not be read as a claim that
NAPQES has been proved DAE-secure. We expect the stronger
indistinguishability-from-random-bits statement to be reachable for this
construction at the cost of the same $q^2/2^{128}$ nonce-collision term the
bounds below already carry, by a union bound over queries in place of the
single-challenge hybrid; establishing it is left to future work.
\end{remark}

\begin{lemma}[Equal padded length is derived]\label{rem:equal-padded-derived}
$|m_0|=|m_1|$ forces $m_0,m_1$ to share the same padding bucket $B$ and the
same real-token count $R$ (Section~\ref{sec:padding}), with certainty —
this is a direct consequence of $B$ depending only on codepoint count.
\end{lemma}

\begin{lemma}[Nonce collision]\label{lem:nonce-collision}
For any two distinct queried pairs $(A_1,M_1)\neq(A_2,M_2)$ under the same
$\mathsf{sk}$, there is a PRF adversary $\mathcal{B}$ with similar running
time such that
$\Pr[N(\mathsf{sk},A_1,M_1)=N(\mathsf{sk},A_2,M_2)] \le \mathrm{Adv}^{\mathrm{PRF}}_{\mathrm{HMAC\text{-}SHA256}}(\mathcal{B}) + 2^{-128}$.
\end{lemma}

\begin{proof}
By Lemma~\ref{lem:domsep}, $(A_1,M_1)\neq(A_2,M_2)$ implies the two domain
$\texttt{0x0A}$ inputs are distinct, so replacing $\mathrm{HMAC}(\mathsf{sk},\cdot)$
by a truly random function $\rho$ (a PRF hop of advantage
$\mathrm{Adv}^{\mathrm{PRF}}_{\mathrm{HMAC\text{-}SHA256}}(\mathcal{B})$)
makes the two 128-bit truncated outputs independent and uniform, colliding
with probability exactly $2^{-128}$.
\end{proof}

\begin{theorem}[IND-CPA-det]\label{thm:ind-cpa}
For any PPT adversary $\mathcal{A}$ respecting
Definition~\ref{def:ind-cpa-det}'s restrictions and making at most $q$
encryption queries, there exists a PRF adversary $\mathcal{B}_1$ with
similar running time such that
\[
  \mathrm{Adv}^{\mathrm{IND\text{-}CPA\text{-}det}}_{\Pi}(\mathcal{A})
  \;\leq\;
  \mathrm{Adv}^{\mathrm{PRF}}_{\mathrm{HMAC\text{-}SHA256}}(\mathcal{B}_1)
  + \frac{q^2}{2^{128}},
\]
against the standard, uniform-key HMAC-SHA256 PRF assumption, with no
key-guessing or key-entropy term.
\end{theorem}

\begin{proof}
We proceed via a sequence of two hybrid games.

\paragraph{Game $G_0$ (real).} The challenger samples $(\mathbf{k},\mathsf{sk})$
and answers queries via the real $\mathsf{Enc}$, which calls
$F(\mathsf{sk},\cdot)$ for all internal derivations.

\paragraph{Game $G_1$ (PRF replacement).} Identical, except
$F(\mathsf{sk},\cdot)$ is replaced by a truly random function
$\rho:\{0,1\}^*\to\{0,1\}^{256}$, sampled lazily.

\begin{lemma}\label{lem:prf-hop}
There is a PRF adversary $\mathcal{B}_1$, running in time
$\mathcal{O}(T_\mathcal{A})$, such that
$|\Pr[\mathcal{A}\text{ wins }G_0]-\Pr[\mathcal{A}\text{ wins }G_1]|\le\mathrm{Adv}^{\mathrm{PRF}}_{\mathrm{HMAC\text{-}SHA256}}(\mathcal{B}_1)$.
\end{lemma}
\begin{proof}
$\mathcal{B}_1$ is given an oracle $\mathcal{O}$ that is either
$F(\mathsf{sk},\cdot)$ or $\rho(\cdot)$. Because $\mathsf{sk}$ is sampled
independently of $\mathbf{k}$ (Definition~\ref{def:keygen}), $\mathcal{B}_1$
samples its own $\mathbf{k}'$ via the prime-sampling step of
$\mathsf{KeyGen}$ to run the arithmetic layer of $\mathsf{Enc}$ directly,
and forwards every domain-derivation call (including domain
$\texttt{0x0A}$) to $\mathcal{O}$. $\mathcal{B}_1$ never needs to know
$\mathsf{sk}$ (or the challenger's hidden key) to run this simulation: when
$\mathcal{O}=F(\mathsf{sk},\cdot)$ the simulation is distributed
identically to $G_0$; when $\mathcal{O}=\rho$, identically to $G_1$.
\end{proof}

\paragraph{Nonce collision under $G_1$.} Let $\mathrm{Coll}$ be the event
that the challenge nonce $N^*$ collides with some queried nonce
$N_1,\ldots,N_q$. By Definition~\ref{def:ind-cpa-det}'s restriction (ii),
every queried pair $(A_i,M_i)$ is distinct from the challenge pair
$(A^*,m_b)$ regardless of $b$, so Lemma~\ref{lem:nonce-collision} applies
validly to every pair $\{(A_i,M_i),(A^*,m_b)\}$; a union bound over the $q$
queries (together with pairwise collisions among $N_1,\ldots,N_q$) gives
$\Pr[\mathrm{Coll}]\le q^2/2^{128}$ (the $\mathrm{Adv}^{\mathrm{PRF}}_{\mathrm{HMAC\text{-}SHA256}}(\mathcal{B}_1)$
term from Lemma~\ref{lem:nonce-collision} is already accounted for by
$\mathcal{B}_1$ above).

\begin{lemma}[Hiding]\label{lem:hiding}
In $G_1$, conditioned on $\overline{\mathrm{Coll}}$,
$\Pr[\mathcal{A}\text{ wins }G_1\mid\overline{\mathrm{Coll}}]=\tfrac12$.
\end{lemma}

\begin{proof}
Under $\overline{\mathrm{Coll}}$, $N^*\notin\{N_1,\ldots,N_q\}$. Every
challenge-phase derivation (domains $\texttt{0x00}$–$\texttt{0x07}$) uses
an input prefixed by its domain byte and then $N^*$; by
Lemma~\ref{lem:domsep} none of these inputs collides with any query-phase
input (different nonce at the fixed offset) or with any other domain's
input (different domain byte). Since $\rho$ is truly random, all
challenge-phase outputs are therefore fresh, uniform, and independent of
the adversary's view. In particular the keystream
$\mathrm{ks}(N^*)$ (domain $\texttt{0x07}$) is independent of the token
blob $B$ it masks (which is itself assembled entirely from domains
$\texttt{0x00},\texttt{0x01},\texttt{0x02},\texttt{0x04},\texttt{0x05},\texttt{0x06}$),
so the one-time-pad identity applies: for fixed $B$,
$\widetilde{B}=B\oplus\mathrm{ks}(N^*)$ is uniform over
$\{0,1\}^{8|B|}$. By Lemma~\ref{rem:equal-padded-derived}, $m_0,m_1$ share
the same padding bucket and real-token count $R$; by
Lemma~\ref{lem:length-det}, $N_{\mathrm{tok}}=R(\mathrm{MAX\_NOISE\_RUN}+1)$
is therefore identical for $m_0$ and $m_1$ (indeed for \emph{any} nonce),
so $|B|$ — and hence the distribution of $\widetilde{B}$ — is identical
regardless of $b$. The tag $T=\rho(\texttt{0x03}\|N^*\|\cdots)$ is a fresh
uniform 32-byte value, also independent of $b$. Hence
$c^*=(N^*,\widetilde{B},T)$ is jointly uniform over its domain,
independently of $b$, giving the adversary no information about $b$.
\end{proof}

Combining: $\Pr[\mathcal{A}\text{ wins }G_1]\le\Pr[\mathrm{Coll}]+\tfrac12\le q^2/2^{128}+\tfrac12$,
so
\[
  \mathrm{Adv}^{\mathrm{IND\text{-}CPA\text{-}det}}_{\Pi}(\mathcal{A})
  = \bigl|\Pr[\mathcal{A}\text{ wins }G_0]-\tfrac12\bigr|
  \le \mathrm{Adv}^{\mathrm{PRF}}_{\mathrm{HMAC\text{-}SHA256}}(\mathcal{B}_1) + \frac{q^2}{2^{128}}. \qedhere
\]
\end{proof}

\subsection{Length Hiding}
\label{sec:length-hiding}

Definition~\ref{def:ind-cpa-det} quantifies only over challenge pairs of
\emph{equal} codepoint length, and therefore says nothing whatsoever about
what an observer learns from $|C|$ itself. The notion below is the one
under which a length-observing adversary must be analysed. It is strictly
stronger than Definition~\ref{def:ind-cpa-det} — the challenge messages may
differ in length — and it is the only sense in which this paper claims
resistance to traffic analysis. Every such claim elsewhere in this paper
refers to Theorem~\ref{thm:lh-ind-cpa} and
Corollary~\ref{cor:length-leak}, and to nothing else.

\begin{definition}[LH-IND-CPA-det security]\label{def:lh-ind-cpa-det}
Fix a padding profile $B(\cdot)$ (Section~\ref{sec:pad-profiles}). The
experiment
$\mathbf{Exp}^{\mathrm{lh\text{-}ind\text{-}cpa\text{-}det}}_{\mathcal{A}}(\lambda)$
is defined exactly as
$\mathbf{Exp}^{\mathrm{ind\text{-}cpa\text{-}det}}_{\mathcal{A}}(\lambda)$
of Definition~\ref{def:ind-cpa-det}, retaining both restrictions (i) and
(ii), except that step 2's condition $|m_0|=|m_1|$ is replaced by the
strictly weaker condition
\[
  B(|m_0|) \;=\; B(|m_1|);
\]
that is, $\mathcal{A}$ may choose challenge messages of \emph{different}
codepoint lengths, subject only to their falling in the same padding
bucket. The advantage
$\mathrm{Adv}^{\mathrm{LH\text{-}IND\text{-}CPA\text{-}det}}_{\Pi}(\mathcal{A})$
is defined as before.
\end{definition}

\begin{theorem}[LH-IND-CPA-det]\label{thm:lh-ind-cpa}
For any PPT adversary $\mathcal{A}$ respecting
Definition~\ref{def:lh-ind-cpa-det}'s restrictions and making at most $q$
encryption queries, there exists a PRF adversary $\mathcal{B}_1$ with
similar running time such that
\[
  \mathrm{Adv}^{\mathrm{LH\text{-}IND\text{-}CPA\text{-}det}}_{\Pi}(\mathcal{A})
  \;\leq\;
  \mathrm{Adv}^{\mathrm{PRF}}_{\mathrm{HMAC\text{-}SHA256}}(\mathcal{B}_1)
  + \frac{q^2}{2^{128}},
\]
i.e.\ exactly the bound of Theorem~\ref{thm:ind-cpa}, with no additional
term for the relaxed length condition.
\end{theorem}

\begin{proof}
The proof of Theorem~\ref{thm:ind-cpa} applies verbatim. Inspecting it, the
hypothesis $|m_0|=|m_1|$ is invoked in exactly one place: the appeal to
Lemma~\ref{rem:equal-padded-derived} inside the proof of
Lemma~\ref{lem:hiding}, whose sole function is to establish that $m_0$ and
$m_1$ share a padding bucket $B$, and hence a real-token count $R=B+2$.
Definition~\ref{def:lh-ind-cpa-det} supplies that conclusion directly as a
hypothesis, so Lemma~\ref{rem:equal-padded-derived} is simply not needed.
Every other step — the PRF hop (Lemma~\ref{lem:prf-hop}), the
nonce-collision bound (Lemma~\ref{lem:nonce-collision}), and the
one-time-pad argument — is independent of the challenge messages' lengths.
By Lemma~\ref{lem:length-det}, $|C|$ is then identical for $m_0$ and $m_1$,
so the ciphertext length carries no information about $b$, and
$\widetilde{B}$ is uniform over strings of that one common length.
\end{proof}

The notion is not vacuous, and it is not satisfied by the schemes NAPQES is
compared against in Section~\ref{sec:comparison}.

\begin{proposition}[Separation from AES-GCM and ChaCha20-Poly1305]
\label{prop:lh-separation}
Let $\Sigma$ be any AEAD scheme whose ciphertext length is a strictly
increasing function of the plaintext length — in particular AES-GCM and
ChaCha20-Poly1305, for which $|C|=|M|+16$ bytes. Then there is an adversary
$\mathcal{A}_{\mathrm{len}}$ making \emph{zero} oracle queries and
performing a single integer comparison such that
\[
  \mathrm{Adv}^{\mathrm{LH\text{-}IND\text{-}CPA\text{-}det}}_{\Sigma}(\mathcal{A}_{\mathrm{len}}) \;=\; \tfrac12,
\]
the maximum possible advantage, irrespective of the strength of the
underlying primitive.
\end{proposition}

\begin{proof}
$\mathcal{A}_{\mathrm{len}}$ outputs any $(A^*,m_0,m_1)$ with
$B(|m_0|)=B(|m_1|)$ but $|m_0|\neq|m_1|$ — for instance $|m_0|=1$ and
$|m_1|=8$, both in bucket $B=16$ under the default profile. On receiving
$c^*$ it returns $b'=0$ if $|c^*|=|m_0|+16$ and $b'=1$ otherwise, winning
with probability $1$. Restrictions (i) and (ii) are satisfied vacuously,
since no oracle query is made.
\end{proof}

\begin{corollary}[Quantified length leakage]\label{cor:length-leak}
Fix a padding profile (Section~\ref{sec:pad-profiles}) with $\beta$
reachable block sizes, and let the plaintext codepoint count $n$ be drawn
from an arbitrary distribution on the message space. Then
\[
  I\bigl(n;\,|C|\bigr) \;\le\; \log_2\beta,
\]
and for a sequence of $m$ ciphertexts,
$I\bigl(n_1,\ldots,n_m;\,|C_1|,\ldots,|C_m|\bigr)\le m\log_2\beta$, with no
independence assumption on the $n_i$. Under the default $\mathsf{bucket}$
profile $\beta=13$, giving at most $\log_2 13\approx3.70$ bits per message;
under $\mathsf{frame}(F)$, $\beta=1$ and the leakage is exactly zero. For
AES-GCM and ChaCha20-Poly1305 on the same message space, $|C|=|M|+16$
determines $n$ exactly, so $I(n;|C|)=H(n)$ — up to $16$ bits per message.
\end{corollary}

\begin{proof}
By Lemma~\ref{lem:length-det}, $|C|$ is a function of $B(n)$, which takes at
most $\beta$ distinct values; hence $H(|C|)\le\log_2\beta$ and
$I(n;|C|)\le H(|C|)\le\log_2\beta$. The $m$-message bound follows from
subadditivity, $H(|C_1|,\ldots,|C_m|)\le\sum_i H(|C_i|)$.
\end{proof}

The next proposition is stated because it is the question a reader
evaluating the cost of this construction will ask, and because its answer
determines where the length-hiding property above actually comes from.

\begin{proposition}[Expansion is length-neutral]\label{prop:expansion-neutral}
Let the wire format induce a map $f$ from padding buckets to ciphertext
lengths, so that $|C|=f(B)$, and suppose $f$ is public and injective. Then
for every prior on the message space and every observer, the posterior
distribution of the plaintext given $|C|$ is the same for every such $f$.
In particular, replacing the token ceiling
$\mathrm{MAX\_NOISE\_RUN}+1=20$ and the $8$-byte token width by \emph{any}
other constants — including $1$ and $1$, that is, deleting the arithmetic
layer outright — leaves the leakage of Corollary~\ref{cor:length-leak}
exactly unchanged.
\end{proposition}

\begin{proof}
If $f$ is public and injective then $B$ and $f(B)$ are computable from one
another by the observer, so the two generate the same $\sigma$-algebra and
condition identically; the claim is then the data-processing equality for a
bijection. The wire format of Section~\ref{sec:wire-format} has
$f(B)=48+8(B+2)(\mathrm{MAX\_NOISE\_RUN}+1)$, strictly increasing in $B$
and hence injective.
\end{proof}

\begin{remark}[Coarsening, not inflation, is what hides length]
\label{rem:coarsening-not-inflation}
Proposition~\ref{prop:expansion-neutral} isolates the mechanism.
Theorem~\ref{thm:lh-ind-cpa} and Corollary~\ref{cor:length-leak} are
consequences of the \emph{many-to-one} padding map $n\mapsto B(n)$ of
Section~\ref{sec:padding}, which collapses $2^{16}$ possible codepoint
counts onto $13$ buckets. They are \emph{not} consequences of the
multi-prime token map, of the noise schedule, or of the fixed
$20$-tokens-per-real-token ceiling: each of those is injective, publicly
invertible post-processing of
$B$, and by Proposition~\ref{prop:expansion-neutral} contributes exactly
zero bits of length confidentiality. It follows that a reader whose sole
objective is traffic-analysis resistance should tune the coarseness of
$B(\cdot)$ (Section~\ref{sec:pad-profiles}) and not the expansion factor —
$\mathsf{frame}(F)$ reduces the leakage to zero, which no choice of
expansion factor can do. Section~\ref{sec:arith-layer-value} states what
the arithmetic layer is retained for instead, and
Section~\ref{sec:lean-profile} gives the construction with that layer
removed, for comparison.
\end{remark}

\paragraph{Measured leakage.}
Corollary~\ref{cor:length-leak} is an upper bound; the realised leakage
depends on the length distribution, so we measure it. The harness
\texttt{traffic\_analysis\_bench.py} in the reference distribution encrypts a
grid of plaintext lengths under each profile and each comparison scheme and
reports, for a uniform prior over $0$--$512$ codepoints on a stride-$8$ grid
($H(n)=6.02$ bits), the realised $I(n;|C|)$, the success rate of a
maximum-a-posteriori adversary recovering $n$ exactly from $|C|$, and the
success rate of a two-class adversary distinguishing a short command
($12$--$20$ codepoints) from a full parameter set ($300$--$500$), where
$50\%$ is a coin flip.

\begin{table}[h]
\centering
\begin{tabular}{@{}lrrrr@{}}
\toprule
Scheme & Distinct $|C|$ & $I(n;|C|)$ & Recover $n$ & Two-class \\
\midrule
NAPQES, $\mathsf{bucket}$ (default) & $7$ & $2.02$ b & $10.8\%$ & $100\%$ \\
NAPQES, $\mathsf{coarse}(3)$        & $3$ & $0.94$ b & $4.6\%$ & $100\%$ \\
NAPQES, $\mathsf{frame}(1024)$      & $1$ & $0.00$ b & $1.5\%$ & $50.0\%$ \\
AES-GCM                             & $65$ & $6.02$ b & $100\%$ & $100\%$ \\
ChaCha20-Poly1305                   & $65$ & $6.02$ b & $100\%$ & $100\%$ \\
\bottomrule
\end{tabular}
\caption{Measured length leakage over a uniform prior on $0$--$512$
codepoints. The two-class column is reported separately because it is not
governed by the average-case bound: a coarse profile can leave $I(n;|C|)$
small while still separating two classes perfectly.}
\label{tab:length-leak-measured}
\end{table}

Two observations from Table~\ref{tab:length-leak-measured} are worth stating
explicitly, since neither is visible from
Corollary~\ref{cor:length-leak} alone. First, the bound is not tight: the
realised $2.02$ bits under the default profile is well below $\log_2 13$,
because a bounded length prior reaches only some of the buckets. Second, and
more importantly, a \emph{small} $I(n;|C|)$ does not imply resistance to a
\emph{specific} distinguisher. Under both $\mathsf{bucket}$ and
$\mathsf{coarse}(3)$ the two-class adversary succeeds with certainty, since
the two classes fall in different buckets ($16$ and $512$) whatever the
expansion factor; only $\mathsf{frame}(F)$ with $F$ above the longer class
reduces it to a coin flip. Deployments that need a particular pair of
message populations to be indistinguishable must therefore choose $F$
larger than both, and must not infer that result from the average-case
bound.

\subsection{INT-CTXT: Ciphertext Integrity}
\label{sec:int-ctxt}

\begin{definition}[INT-CTXT security]\label{def:int-ctxt}
For a PPT adversary $\mathcal{A}$, define
$\mathbf{Exp}^{\mathrm{int\text{-}ctxt}}_{\mathcal{A}}(\lambda)$:
\begin{enumerate}
  \item $(\mathbf{k},\mathsf{sk})\gets\mathsf{KeyGen}(1^\lambda)$, and the
        challenger initialises an empty set $\mathcal{S}$.
  \item $\mathcal{A}$ is given adaptive access to two oracles, interleaved
        in any order of its choosing:
        \begin{itemize}
          \item $\mathcal{O}_{\mathsf{Enc}}(A,M)$ returns
                $c\gets\mathsf{Enc}(\mathbf{k},\mathsf{sk},A,M)$ and adds
                the pair $(c,A)$ to $\mathcal{S}$. At most $q$ such queries
                are made.
          \item $\mathcal{O}_{\mathsf{Vf}}(A,c)$ computes
                $\mathsf{Dec}(\mathbf{k},\mathsf{sk},A,c)$ and returns the
                single bit $[\mathsf{Dec}(\cdot)\neq\bot]$ — the
                verification oracle reveals \emph{only} acceptance or
                rejection, never the recovered plaintext. At most $q_v\ge1$
                such queries are made.
        \end{itemize}
  \item The experiment outputs $1$ iff at some point $\mathcal{A}$ submits
        a \emph{fresh} query $\mathcal{O}_{\mathsf{Vf}}(A^*,c^*)$ — one with
        $(c^*,A^*)\notin\mathcal{S}$ \emph{at the moment of submission} —
        that returns $1$.
\end{enumerate}
The experiment does not halt on a successful forgery: $\mathcal{A}$ may
continue querying both oracles, and wins if \emph{any} of its $q_v$
verification queries is fresh and accepting. The advantage of $\mathcal{A}$
is
$\mathrm{Adv}^{\mathrm{INT\text{-}CTXT}}_{\Pi}(\mathcal{A}) =
\Pr[\mathbf{Exp}^{\mathrm{int\text{-}ctxt}}_{\mathcal{A}}(\lambda)=1]$.
Freshness is measured over the \emph{pair} $(c,A)$, so a replay of a
previously returned ciphertext under a different associated-data string
counts as a forgery attempt, as it must for an AEAD scheme.
\end{definition}

\begin{theorem}[INT-CTXT]\label{thm:int-ctxt}
Let $\mathcal{A}$ be any PPT adversary in the experiment of
Definition~\ref{def:int-ctxt}, making at most $q$ encryption queries
and at most $q_v\ge1$ adaptive verification queries. There
exists a PRF adversary $\mathcal{B}_2$ with similar running time such that
\[
  \mathrm{Adv}^{\mathrm{INT\text{-}CTXT}}_{\Pi}(\mathcal{A})
  \leq
  \mathrm{Adv}^{\mathrm{PRF}}_{\mathrm{HMAC\text{-}SHA256}}(\mathcal{B}_2)
  + \frac{q_v}{2^{256}}.
\]
There is no $q$-dependent (encryption-query-count) term, since Case~2
below is resolved deterministically rather than probabilistically.
\end{theorem}

\begin{proof}
Write a forgery attempt as $c^*=(N^*,\widetilde{B}^*,T^*)$ with tag input
$x^*=\texttt{0x03}\|N^*\|\mathrm{TagIn}(aad^*,\widetilde{B}^*)$
(Definition~\ref{def:taginput}), the same quantity $\mathsf{Dec}$ step (4)
recomputes.

\paragraph{Case 1: $N^*$ was not used in any encryption query.}
Lemma~\ref{lem:domsep} is a statement about the injectivity of the input
encoding over its whole argument domain, so it applies to $x^*$ even
though $\mathcal{A}$, not the construction, chose the fields of $x^*$.
By that lemma, $x^*$ is distinct from every query-phase domain
$\texttt{0x03}$ input (different nonce at the fixed offset) and from every
other domain's input, so $x^*$ was never submitted to
$\mathrm{HMAC}(\mathsf{sk},\cdot)$. Construct $\mathcal{B}_2$, given an
oracle $f$ (real or random), exactly as $\mathcal{B}_1$ in
Lemma~\ref{lem:prf-hop}: sample $\mathbf{k}'$ independently, forward every
HMAC call to $f$, and output ``real'' iff some forgery attempt's candidate
tag matches $f(x^*)$. In the ideal world $f(x^*)$ is uniform and
independent of $\mathcal{A}$'s view, so each of the $q_v$ attempts succeeds
by blind guessing with probability exactly $2^{-256}$; a union bound gives
$\Pr[\text{forge in Case 1}]\le\mathrm{Adv}^{\mathrm{PRF}}_{\mathrm{HMAC\text{-}SHA256}}(\mathcal{B}_2)+q_v/2^{256}$.

\paragraph{Case 2: $N^*=N_i$ for some query $i$.} If $x^*\neq x_i$ for
every such $i$ (which, by injectivity of the length-prefixed AAD encoding,
covers both $\widetilde{B}^*\neq\widetilde{B}_i$ and the
$\widetilde{B}^*=\widetilde{B}_i\wedge aad^*\neq aad_i$ replay-with-new-AAD
sub-case), $x^*$ was never queried and Case 1's argument applies verbatim.
If instead $x^*=x_i$ for some $i$, injectivity forces
$aad^*=aad_i\wedge\widetilde{B}^*=\widetilde{B}_i$: if also $T^*=T_i$ then
$(c^*,aad^*)=(c_i,aad_i)$, excluded by freshness; if $T^*\neq T_i$, the
recomputed tag equals $T_i\neq T^*$ and $\mathsf{Dec}$ rejects
unconditionally. Either sub-case contributes no successful,
freshness-respecting forgery, and both branches are resolved
deterministically via Lemma~\ref{lem:domsep}'s injectivity — $N^*$ is
chosen adaptively by $\mathcal{A}$ after observing $N_1,\ldots,N_q$, so no
probabilistic collision term applies to Case~2.

Combining both cases yields the stated bound.
\end{proof}

\subsection{IND-CCA-det Security}
\label{sec:ind-cca}

\begin{definition}[IND-CCA-det security]\label{def:ind-cca-det}
For a PPT adversary $\mathcal{A}=(\mathcal{A}_1,\mathcal{A}_2)$, define
$\mathbf{Exp}^{\mathrm{ind\text{-}cca\text{-}det}}_{\mathcal{A}}(\lambda)$
exactly as
$\mathbf{Exp}^{\mathrm{ind\text{-}cpa\text{-}det}}_{\mathcal{A}}(\lambda)$
of Definition~\ref{def:ind-cpa-det}, with these additions:
\begin{itemize}
  \item In both phases $\mathcal{A}$ additionally has adaptive access to a
        decryption oracle
        $\mathcal{O}_{\mathsf{Dec}}(A,c)=\mathsf{Dec}(\mathbf{k},\mathsf{sk},A,c)$,
        returning the recovered plaintext or $\bot$, making at most $q_d$
        such queries in total.
  \item After the challenge $c^*$ on associated data $A^*$ has been issued,
        $\mathcal{A}_2$ may query $\mathcal{O}_{\mathsf{Dec}}$ on anything
        except the exact pair $(A^*,c^*)$. In particular $(A,c^*)$ for
        $A\neq A^*$ is permitted, as in the standard AEAD formulation.
\end{itemize}
Restrictions (i) and (ii) of Definition~\ref{def:ind-cpa-det} bind
$\mathcal{A}$'s \emph{encryption}-oracle queries in both phases, and are
not imposed on decryption queries. The advantage of $\mathcal{A}$ is
$\mathrm{Adv}^{\mathrm{IND\text{-}CCA\text{-}det}}_{\Pi}(\mathcal{A}) =
|\Pr[\mathbf{Exp}^{\mathrm{ind\text{-}cca\text{-}det}}_{\mathcal{A}}(\lambda)=1]-\tfrac12|$.
\end{definition}

\begin{theorem}[IND-CCA-det]\label{thm:ind-cca}
For any PPT adversary $\mathcal{A}$ in the experiment of
Definition~\ref{def:ind-cca-det}, making at most $q$ encryption queries and
$q_d$ decryption queries, there exist
PRF adversaries $\mathcal{B}_1,\mathcal{B}_2$ with similar running time
such that
\[
  \mathrm{Adv}^{\mathrm{IND\text{-}CCA\text{-}det}}_{\Pi}(\mathcal{A})
  \leq
  \mathrm{Adv}^{\mathrm{PRF}}_{\mathrm{HMAC\text{-}SHA256}}(\mathcal{B}_1)
  +
  \mathrm{Adv}^{\mathrm{PRF}}_{\mathrm{HMAC\text{-}SHA256}}(\mathcal{B}_2)
  + \frac{q^2}{2^{128}}
  + \frac{q_d}{2^{256}}.
\]
\end{theorem}

\begin{proof}
The argument is the standard encrypt-then-authenticate composition of
Bellare and Namprempre~\cite{bellare2000ae}, but that result is stated for
the unrestricted IND-CPA notion, which no deterministic scheme satisfies
(Section~\ref{sec:ind-cpa}). We therefore do not invoke it as a black box;
we give the two game hops in full against
Definitions~\ref{def:ind-cpa-det}, \ref{def:int-ctxt}
and~\ref{def:ind-cca-det}. Nothing beyond the restricted notions is used.

\paragraph{Game $H_0$ (real).} The IND-CCA-det experiment of
Definition~\ref{def:ind-cca-det}, verbatim.

\paragraph{Game $H_1$ (table-backed decryption).} As $H_0$, except that the
challenger maintains a table $T$ recording $(c,A)\mapsto M$ for every
encryption-oracle answer $c$ on input $(A,M)$, together with
$(c^*,A^*)\mapsto m_b$ once the challenge is issued, and answers a
decryption query $(A,c)$ with $T[(c,A)]$ if that entry exists and with
$\bot$ otherwise.

\paragraph{Transition.} Let $D$ be the event that $\mathcal{A}$ submits a
decryption query $(A,c)$ with $(c,A)\notin T$ for which
$\mathsf{Dec}(\mathbf{k},\mathsf{sk},A,c)\neq\bot$. On every query with
$(c,A)\in T$ the two games agree: $H_1$ returns the recorded plaintext, and
$H_0$ returns the same value by correctness
(Proposition~\ref{prop:correctness}). On every query with $(c,A)\notin T$
the two games agree unless $D$ occurs. Hence $H_0$ and $H_1$ are identical
until $D$, and $|\Pr[H_0{=}1]-\Pr[H_1{=}1]|\le\Pr[D]$.

\paragraph{Bounding $\Pr[D]$.} Construct an INT-CTXT adversary
$\mathcal{A}'$ (Definition~\ref{def:int-ctxt}) that runs $\mathcal{A}$,
maintaining its own copy of $T$. It forwards each of $\mathcal{A}$'s $q$
encryption queries to $\mathcal{O}_{\mathsf{Enc}}$, recording the answers in
$T$; it samples $b$ itself and obtains the challenge $c^*$ by a
$(q{+}1)$-th query $\mathcal{O}_{\mathsf{Enc}}(A^*,m_b)$, recording
$(c^*,A^*)\mapsto m_b$. A decryption query $(A,c)$ with $(c,A)\in T$ is
answered from $T$; one with $(c,A)\notin T$ is forwarded to
$\mathcal{O}_{\mathsf{Vf}}(A,c)$, and $\mathcal{A}'$ answers $\bot$ to
$\mathcal{A}$ if the oracle returns $0$, halting immediately if it returns
$1$. Note that $\mathcal{A}'$ only ever needs the verification oracle's
single bit, never a decrypted plaintext, so the simulation is compatible
with the bit-only oracle of Definition~\ref{def:int-ctxt}. Every query
$\mathcal{A}'$ forwards to $\mathcal{O}_{\mathsf{Vf}}$ is fresh, because
$T$ contains exactly the pairs the encryption oracle returned. The
simulation is perfect up to and including the first occurrence of $D$,
which is all that is required, and $D$ occurring is precisely
$\mathcal{A}'$ winning. With $q_v=q_d$ verification queries,
Theorem~\ref{thm:int-ctxt} gives
$\Pr[D]\le\mathrm{Adv}^{\mathrm{INT\text{-}CTXT}}_{\Pi}(\mathcal{A}')
\le\mathrm{Adv}^{\mathrm{PRF}}_{\mathrm{HMAC\text{-}SHA256}}(\mathcal{B}_2)+q_d/2^{256}$.
The extra $(q{+}1)$-th encryption query costs nothing, since that bound
carries no encryption-query-count term.

\paragraph{Bounding $|\Pr[H_1{=}1]-\tfrac12|$.} Construct an IND-CPA-det
adversary $\mathcal{A}''$ (Definition~\ref{def:ind-cpa-det}) that runs
$\mathcal{A}$, forwards its encryption queries to its own encryption
oracle while recording $T$, forwards $(A^*,m_0,m_1)$ as its own challenge
and relays the challenger's $c^*$, and answers decryption queries from $T$,
returning $\bot$ for entries it does not hold. This is exactly $H_1$'s
oracle: the only table entry $\mathcal{A}''$ cannot populate is
$(c^*,A^*)\mapsto m_b$, and that entry is never read, since
Definition~\ref{def:ind-cca-det} forbids $\mathcal{A}$ the pair
$(A^*,c^*)$. Any permitted query $(A,c^*)$ with $A\neq A^*$ is a distinct
table key, absent from $T$ in $H_1$ as well, so both answer $\bot$.
$\mathcal{A}''$ inherits restrictions (i) and (ii) from $\mathcal{A}$ and
makes the same $q$ encryption queries, so by
Theorem~\ref{thm:ind-cpa},
$|\Pr[H_1{=}1]-\tfrac12|
\le\mathrm{Adv}^{\mathrm{PRF}}_{\mathrm{HMAC\text{-}SHA256}}(\mathcal{B}_1)+q^2/2^{128}$.

Combining the two bounds yields the stated inequality.
\end{proof}

% ─────────────────────────────────────────────────────────────────────────────
\section{Post-Quantum Analysis}
\label{sec:pq}
% ─────────────────────────────────────────────────────────────────────────────

\paragraph{SHA-256 and Grover's algorithm.}
Grover's algorithm~\cite{grover96} provides a quadratic speedup for
unstructured search, requiring $\approx2^{n/2}$ oracle queries against an
$n$-bit target. Applied to SHA-256 preimage search, this reduces the
nominal security level from 256 to approximately 128 bits. This $2^{128}$
figure is a bound on \emph{sequential query depth}, not a parallelizable
work-factor: the Bennett–Bernstein–Brassard–Vazirani (BBBV) optimality
result~\cite{bbbv1997} shows no quantum algorithm can search an
unstructured space of size $N$ in fewer than $\Omega(\sqrt N)$ oracle
queries, and Zalka~\cite{zalka1999} shows that distributing Grover's
search across $S$ machines yields only a $\sqrt{S}$ speedup — the depth
stays $\approx2^{n/2}/\sqrt{S}$ and the total work is unchanged, unlike
classical brute force, which parallelizes linearly. NIST's transition
guidance~\cite{sp800-131a} deprecates 112-bit security after 2030 and sets
128 bits as the minimum from that point; at $2^{128}$ sequential oracle
queries this construction meets that post-2030 minimum, rather than merely
exceeding the 112-bit level being retired.

\paragraph{No Shor-applicable structure.}
NAPQES does not use discrete logarithms, integer factorisation, or any
algebraic group structure, so Shor's algorithm provides no speedup against
it. This is a property NAPQES \emph{shares} with AES-GCM,
ChaCha20-Poly1305 and Ascon, all of which are likewise symmetric and
likewise unaffected by Shor; it is recorded here to forestall the common
misreading that a symmetric scheme's arithmetic content makes it
quantum-vulnerable, and it is not offered as a differentiator. The
differentiator claimed in this paper is the absence of a polynomial
authenticator and the resulting absence of the nonce-reuse key-recovery
mode of Section~\ref{sec:related}.

\paragraph{Two separate key-space figures.}
Because $\mathsf{sk}$ is the sole HMAC key and $\mathbf{k}$ is used only by
the un-keyed arithmetic layer (Remark~\ref{rem:key-roles}), the two secrets
have different, non-interchangeable security roles and are best reported
separately rather than as one conflated figure:
\begin{itemize}
  \item $\mathsf{sk}$: uniform 256 bits, classical exhaustive search
        $2^{256}$, post-Grover sequential depth $\approx2^{128}$ — this is
        the figure relevant to Assumption~\ref{assum:hmac-prf} and to every
        theorem of Section~\ref{sec:security}.
  \item $\mathbf{k}$: $\approx2^{191.46}$ ordered-tuple key-space size at the
        default $K=10$ (post-Grover $\approx2^{95.73}$). This figure is
        reported for completeness only. It is \emph{not} relevant to any
        theorem in Section~\ref{sec:security}, none of which depend on
        $H_\infty(\mathbf{k})$ (Remark~\ref{rem:key-roles}), nor to the
        length-hiding property of Section~\ref{sec:length-hiding}, which is
        independent of the arithmetic layer entirely
        (Remark~\ref{rem:coarsening-not-inflation}). What that layer is
        retained for is set out in Section~\ref{sec:arith-layer-value}.
\end{itemize}
Both figures are non-parallelizable sequential-depth bounds by the same
BBBV~\cite{bbbv1997}/Zalka~\cite{zalka1999} argument. Section~\ref{sec:caveats}
discusses the practical consequence of this asymmetry: protecting
$\mathsf{sk}$ is necessary and sufficient for confidentiality and
integrity; protecting $\mathbf{k}$ in addition provides no independent
confidentiality margin once $\mathsf{sk}$ is compromised.

% ─────────────────────────────────────────────────────────────────────────────
\section{What the Arithmetic Layer Does and Does Not Buy}
\label{sec:arith-layer-value}
% ─────────────────────────────────────────────────────────────────────────────

The arithmetic layer — the multi-prime token map $c\mapsto c\cdot k_j+a$,
the noise schedule, and the $20$-token ceiling — accounts for essentially
the whole cost of this construction, and for none of its theorems. A reader
deciding whether to deploy NAPQES is entitled to a direct statement of what
that cost purchases. We give one, separated into what is proved, what is
argued but not proved, and what is not claimed at all.

\paragraph{Proved: the noise layer is length-neutral.}
Lemma~\ref{lem:length-det} is a property \emph{of} the arithmetic layer,
but a defensive one: the $\mathrm{MAX\_NOISE\_RUN}$ cap together with the
ceiling-padding step exists to ensure that the noise schedule — which is
derived from the message-dependent synthetic nonce
(Definition~\ref{def:synthnonce}) — does not itself introduce a
message-dependent length. Without the cap and ceiling, the natural token
count would vary with the message even at a fixed bucket, and an observer
collecting several encryptions of one message under varying AAD could
average that variance away and recover the bucket, and more. The cap and
ceiling remove that channel entirely. This is a genuine contribution of the
layer's design, but it is the removal of a leak that the layer would
otherwise create; it is not an improvement on a construction that has no
noise layer to begin with.

\paragraph{Argued, not proved: degradation under nonce collision.}
Lemma~\ref{lem:nonce-collision} bounds the probability of two distinct
$(A,M)$ pairs sharing a nonce by $\mathrm{Adv}^{\mathrm{PRF}}+2^{-128}$, so
what follows lies outside the model of Section~\ref{sec:security} and is
\emph{not} claimed as a security property. Should such a collision
nevertheless occur, the two ciphertexts share a keystream and the adversary
obtains $B_1\oplus B_2$. In a construction that masks padded codepoints
directly, this is a textbook two-time pad on the plaintexts, recoverable by
standard crib-dragging. Here it is the XOR of two token blobs: the noise
and filler positions, being derived from the shared nonce alone, cancel to
zero and are thereby revealed, while each real position yields the XOR of
two affine images $c\cdot k_j+a$ and $c'\cdot k_j+a$ sharing an unknown
multiplier $k_j$ and a common addend. Recovering codepoints therefore
requires additionally resolving $\mathbf{k}$. We record this as strictly
more work for the adversary than a plain two-time pad, and as nothing more:
it is substantial plaintext leakage in both cases, it admits no bound we
are able to state, and it must not be read as a defence.

\paragraph{Not claimed.}
Three claims that this layer might be expected to support are not made
anywhere in this paper, and we state their absence explicitly.
\begin{itemize}
  \item \emph{No content confidentiality.} The one-time-pad step of
        Lemma~\ref{lem:hiding} is established from domain $\texttt{0x07}$
        alone and never inspects the token blob it masks. The arithmetic
        layer contributes nothing to Theorem~\ref{thm:ind-cpa},
        Theorem~\ref{thm:int-ctxt}, or Theorem~\ref{thm:ind-cca}, and
        removing it changes none of their statements or bounds
        (Proposition~\ref{prop:lean-theorems}).
  \item \emph{No length hiding.} By
        Proposition~\ref{prop:expansion-neutral} the expansion factor is a
        public injective function of the padding bucket and contributes
        exactly zero bits of length confidentiality; all of the property
        proved in Section~\ref{sec:length-hiding} comes from the coarseness
        of $B(\cdot)$, which survives the removal of the arithmetic layer
        untouched (Remark~\ref{rem:coarsening-not-inflation}).
  \item \emph{No dependence on $K$ or on $H_\infty(\mathbf{k})$.} $K$ is an
        interoperability parameter, not a security parameter
        (Remark~\ref{rem:key-roles}), and $\mathbf{k}$ adds no
        confidentiality margin once $\mathsf{sk}$ is known
        (Section~\ref{sec:caveats}).
\end{itemize}

\subsection{The Construction Without the Arithmetic Layer}
\label{sec:lean-profile}

To make the trade-off precise rather than rhetorical, we specify the same
construction with the arithmetic layer deleted and compare the two. Define
\emph{NAPQES-L} identically to Definition~\ref{def:aead-triple} — same key
generation, same domain separation, same padding and padding profiles, same
synthetic nonce, same domain-$\texttt{0x07}$ keystream, same
domain-$\texttt{0x03}$ tag over the same $\mathrm{TagIn}$ — except that the
token-emission loop of Section~\ref{sec:token-construction} is replaced by
the identity map on padded codepoints, each serialised as $3$ bytes
big-endian. Domains $\texttt{0x00}$, $\texttt{0x01}$, $\texttt{0x02}$,
$\texttt{0x04}$ and $\texttt{0x05}$ are then unused, $\mathbf{k}$ and $K$
disappear from the scheme, and
\[
  |C|_{\text{NAPQES-L}} \;=\; 16 + 3(B+2) + 32 \;=\; 48+3(B+2)\ \text{bytes},
\]
against $|C| = 48+160(B+2)$ for NAPQES as specified; at the minimum bucket
$B=16$ this is $102$ bytes against $2928$.

\begin{proposition}[The theorems are indifferent to the arithmetic layer]
\label{prop:lean-theorems}
Theorem~\ref{thm:ind-cpa}, Theorem~\ref{thm:lh-ind-cpa},
Theorem~\ref{thm:int-ctxt} and Theorem~\ref{thm:ind-cca} hold for NAPQES-L
with identical statements and identical bounds, by the same proofs.
\end{proposition}

\begin{proof}
Each proof uses exactly four facts about the masked payload: that domain
separation holds (Lemma~\ref{lem:domsep}), which is unchanged; that the
nonce-collision bound holds (Lemma~\ref{lem:nonce-collision}), which
depends only on domain $\texttt{0x0A}$ and is unchanged; that the payload
is assembled entirely from domains other than $\texttt{0x07}$, so that the
one-time-pad step of Lemma~\ref{lem:hiding} applies — in NAPQES-L the
payload is the padded codepoint sequence, built from the plaintext and
domain $\texttt{0x06}$ alone, so this still holds; and that the payload
length is a function of $B$ alone, which for NAPQES-L is immediate, since
it is exactly $3(B+2)$ bytes. No proof step examines the token map, the
noise schedule, or $\mathbf{k}$.
\end{proof}

\begin{table}[h]
\centering
\begin{tabular}{@{}lcc@{}}
\toprule
 & NAPQES (specified) & NAPQES-L \\
\midrule
IND-CPA-det, INT-CTXT, IND-CCA-det & identical bounds & identical bounds \\
LH-IND-CPA-det (Thm.~\ref{thm:lh-ind-cpa}) & holds & holds \\
Length leakage (Cor.~\ref{cor:length-leak}) & $\le\log_2\beta$ & $\le\log_2\beta$ \\
Ciphertext size, bucket $B$ & $48+160(B+2)$ B & $48+3(B+2)$ B \\
Minimum ciphertext & $2928$ B & $102$ B \\
Secrets & $\mathbf{k}$ and $\mathsf{sk}$ & $\mathsf{sk}$ only \\
Trial-division caveat (\S\ref{sec:caveats}) & applies & vacuous \\
Nonce-collision residue & affine, $\mathbf{k}$ unknown & two-time pad \\
\bottomrule
\end{tabular}
\caption{NAPQES as specified against the same construction with the
arithmetic layer removed. No theorem of Section~\ref{sec:security} changes.}
\label{tab:lean-comparison}
\end{table}

NAPQES-L is stated for comparison only. It is \textbf{not} part of this
specification, has no wire format, no test vectors and no reference
implementation, and is not interoperable with NAPQES. We include it because
the honest summary of Table~\ref{tab:lean-comparison} is that the
arithmetic layer costs a factor of ${\approx}53$ in bandwidth and buys no
theorem; a deployment selecting NAPQES over NAPQES-L is selecting the
beyond-model margin described above, and should do so deliberately.

% ─────────────────────────────────────────────────────────────────────────────
\section{Comparison with AES-GCM, ChaCha20-Poly1305, and Ascon}
\label{sec:comparison}
% ─────────────────────────────────────────────────────────────────────────────

\begin{table}[h]
\centering
\begin{tabular}{@{}lcccc@{}}
\toprule
Property & NAPQES & AES-GCM & ChaCha20-Poly1305 & Ascon \\
\midrule
Primitive basis & HMAC-SHA256 & AES + GHASH & ChaCha20 + Poly1305 & Ascon permutation \\
FIPS-approved primitives only & \textbf{Yes} & Yes & No & Yes~\cite{sp800-232} \\
Approved as a mode/algorithm & No & Yes & No & Yes~\cite{sp800-232} \\
Shor-applicable structure & No & No & No & No \\
Forbidden Attack risk & \textbf{No} & Yes (nonce reuse) & Yes (nonce reuse) & No \\
Nonce-reuse key recovery & \textbf{No} (Definition~\ref{def:synthnonce}, disclosed message-equality only) & Yes (GCM key) & Yes (Poly1305 key) & Partial \\
Ciphertext expansion & $160\times$ (fixed) & $1\times$ & $1\times$ & $1\times$ \\
Length leakage per message & $\le3.70$ bits ($0$ with $\mathsf{frame}$) & $=H(n)$ & $=H(n)$ & $=H(n)$ \\
Key size & $5K+32$ bytes ($\mathbf{k}$+$\mathsf{sk}$; $82$ at $K=10$) & 32 bytes & 32 bytes & 16 bytes \\
\bottomrule
\end{tabular}
\caption{Comparison of NAPQES with established AEAD schemes. The two FIPS
rows are distinct and must not be conflated: NAPQES uses only an approved
primitive (HMAC-SHA256, FIPS 198-1/180-4), but NAPQES \emph{itself} is not
an approved algorithm or mode. In a FIPS 140-3 validation it would fall
outside the approved-mode boundary and would be reported as a
non-approved algorithm, notwithstanding the provenance of its primitive.}
\label{tab:comparison}
\end{table}

The primary trade-off is ciphertext expansion. The token-emission loop caps
consecutive noise tokens at $\mathrm{MAX\_NOISE\_RUN}=19$
(Section~\ref{sec:token-construction}) and pads every message to exactly
$R\cdot(\mathrm{MAX\_NOISE\_RUN}+1)=20R$ tokens of $8$ bytes each, so the
expansion is a hard, deterministic $160$ bytes per padded codepoint — never
a range, and never dependent on the per-message noise realisation
(Lemma~\ref{lem:length-det}). We do not offer this expansion as a
traffic-analysis countermeasure; by
Proposition~\ref{prop:expansion-neutral} it provides none, and
Section~\ref{sec:lean-profile} gives the same construction without it,
under which every theorem of Section~\ref{sec:security} continues to hold
with identical bounds. What NAPQES does provide against a length-observing
adversary, and AES-GCM and ChaCha20-Poly1305 do not
(Proposition~\ref{prop:lh-separation}), is the length-hiding property of
Section~\ref{sec:length-hiding}, which follows from the padding profile
alone and is tunable down to zero leakage via $\mathsf{frame}(F)$
(Section~\ref{sec:pad-profiles}).

Unlike AES-GCM and ChaCha20-Poly1305, where nonce reuse degrades
confidentiality but does not expose the key, NAPQES's synthetic,
message-bound nonce (Definition~\ref{def:synthnonce}) removes nonce
selection from the caller entirely, so there is no DRBG-failure or
repeated-counter path to reuse; re-querying with an \emph{identical}
$(A,M)$ pair reproduces the identical ciphertext deterministically,
disclosing only plaintext equality — the standard, disclosed MRAE
trade-off (cf.\ AES-GCM-SIV~\cite{rfc8452}), not a confidentiality or
key-recovery break.

This does \emph{not} eliminate the birthday bound. Two distinct $(A,M)$
pairs collide on the truncated 128-bit nonce with probability $2^{-128}$
each (Lemma~\ref{lem:nonce-collision}), and the $q^2/2^{128}$ term in
Theorems~\ref{thm:ind-cpa} and~\ref{thm:ind-cca} is exactly the union bound
over the $\binom{q}{2}$ pairs — an ordinary birthday term. NAPQES therefore
carries a data-complexity limit of roughly $2^{64}$ messages per key, of
the same order as any 128-bit-nonce AEAD, and the practical difference from
AES-GCM is in \emph{how} a nonce collision arises (a PRF collision on
distinct inputs, never a caller or DRBG error) and in what it costs when it
does, not in whether the birthday bound applies. Section~\ref{sec:caveats}
states the recommended operational cap.

% ─────────────────────────────────────────────────────────────────────────────
\section{Implementation}
\label{sec:implementation}
% ─────────────────────────────────────────────────────────────────────────────

\subsection{Python Reference Implementation}

The Python reference implementation (\texttt{napqes.py}) is pure Python
3.10+, with no dependencies beyond the standard library (\texttt{hmac},
\texttt{hashlib}, \texttt{secrets}, \texttt{base64}). Key generation and
the AEAD entry points are \texttt{generate\_v8\_key},
\texttt{encrypt\_bytes\_v8}, and \texttt{decrypt\_bytes\_v8}. All KAT
vectors are derived from this implementation.

\subsection{Rust Implementation}

A Rust port (\texttt{rust/src/lib.rs}) mirrors the Python implementation
(\texttt{generate\_v8\_key}, \texttt{encrypt\_bytes\_v8},
\texttt{decrypt\_bytes\_v8}). Unit tests cover round-trip correctness,
tamper/AAD-mismatch rejection, same-message determinism (the disclosed
MRAE trade-off), distinct-message nonce distinctness, and
ciphertext-length determinism across varied AAD and keys
(Lemma~\ref{lem:length-det}).

\subsection{C Implementation}

A C port (\texttt{C/napqes.c}) mirrors the same construction
(\texttt{napqes\_encrypt\_bytes\_v8}, \texttt{napqes\_decrypt\_bytes\_v8}).

\subsection{Constant-Time Considerations}
\label{sec:constant-time}

\paragraph{Threat model.} Every claim in Section~\ref{sec:security} is made
in the standard black-box model, in which the adversary sees only oracle
inputs and outputs. Side channels — wall-clock time, cache state, branch
predictor state, power, electromagnetic emanation — are explicitly outside
that model, and none of the theorems in this paper says anything about
them. This subsection records what the reference implementations do and do
not do; it is not a claim of side-channel resistance, and the phrase
``compared in constant time'' in $\mathsf{Dec}$ step (5)
(Definition~\ref{def:aead-triple}) should be read as describing exactly one
operation, not the algorithm as a whole.

\paragraph{What is constant-time.} Authentication tag comparison uses
\texttt{hmac.compare\_digest} (Python), \texttt{constant\_time\_eq} from
the \texttt{subtle} crate (Rust), and an OR-accumulator byte-difference
loop (C), so that comparison timing does not depend on the position of the
first differing tag byte. This is the property that matters for forgery
attempts, since it prevents byte-at-a-time tag search. The C accumulator is
not declared \texttt{volatile}, so nothing in the C standard prevents a
sufficiently aggressive compiler from reintroducing an early exit; this is
the usual caveat for hand-rolled constant-time C and is not addressed here.

\paragraph{What is not.} The token-emission and token-inversion loops are
\emph{not} constant-time, and are not intended to be:
\begin{itemize}
  \item The number of loop iterations depends on the realised noise
        pattern, which is a deterministic function of $(\mathsf{sk},N)$.
        The total ciphertext length is padded to a fixed ceiling
        (Lemma~\ref{lem:length-det}) and so leaks nothing, but the
        \emph{time} to produce or consume it varies with the number of
        noise tokens actually emitted before the run cap engages.
  \item Both directions use arbitrary-precision or 64-bit integer
        multiplication and division by a key element $k$. Division latency
        is data-dependent on many microarchitectures.
  \item The padding loop of Section~\ref{sec:padding} performs $B-n$
        derivations, so encryption time depends on the \emph{exact}
        codepoint count $n$, not merely on its padding bucket $B$. The
        padding-bucket caveat below concedes a $\log_2 13\approx3.70$-bit
        length leak through ciphertext \emph{length}; a timing channel can
        in principle recover more than that.
  \item The Python reference implementation offers no timing guarantees of
        any kind: it uses \texttt{int}, list resizing, and dictionary
        lookups throughout.
  \item The structural checks added at $\mathsf{Dec}$ steps (6) and (8)
        (Remark~\ref{rem:dec-structural}) reject at different points and in
        different time. They are reachable only after the tag has verified,
        i.e.\ only by a party already holding $\mathsf{sk}$, so this
        difference is not an oracle for an outside adversary.
\end{itemize}
An attacker with a precise timing channel against a decrypting party may
therefore learn something about the noise realisation, and hence about
$\theta(N)$ and indirectly about $N$; $N$ is transmitted in the clear in
any case, so this particular leak is not by itself a break, but no bound is
offered on what else such a channel exposes. A dudect-style timing study of
the C port is reported in \texttt{docs/DUDECT\_ATTESTATION.md}; it covers
the tag comparison only. Deployments in which the adversary can measure
decryption time on attacker-influenced inputs should treat this as an open
gap and are directed to Section~\ref{sec:caveats}.

% ─────────────────────────────────────────────────────────────────────────────
\section{NIST SP 800-22 Statistical Analysis}
\label{sec:sts}
% ─────────────────────────────────────────────────────────────────────────────

We applied the full NIST SP 800-22 Rev.\ 1a Statistical Test
Suite~\cite{sp800-22} to $10^7$ bits of NAPQES ciphertext output generated
by encrypting a cycling corpus of printable-ASCII plaintext with a
10-element prime key. All 15 eligible SP 800-22 tests passed at the
$\alpha=0.01$ significance level. Results are reproducible using
\texttt{tests/sts\_pipeline.py} in the reference implementation
repository.

% ─────────────────────────────────────────────────────────────────────────────
\section{Known-Answer Tests}
\label{sec:kats}
% ─────────────────────────────────────────────────────────────────────────────

The corpus for the scheme specified in this paper is
\texttt{tests/kat/v8\_vectors.json}. It contains 20 vectors: 12 positive
(ids \texttt{W001}--\texttt{W012}) covering the empty message, a single
character, the block boundaries at 15, 16 and 32 codepoints, 1-element and
10-element prime tuples $\mathbf{k}$, empty, binary and 64-byte associated
data (the last pinning the 8-byte $\mathrm{be8}$ AAD length prefix of
Definition~\ref{def:taginput}), and mixed-case punctuation; and 8 negative.
Of the negative vectors, \texttt{W-N01}--\texttt{W-N05} exercise
authentication (flipped tag byte, flipped masked-blob byte, wrong $A$,
wrong $\mathsf{sk}$) and the 48-byte parse floor, while
\texttt{W-N06}--\texttt{W-N08} are \emph{validly tagged} but structurally
malformed and therefore reach the post-authentication checks of
$\mathsf{Dec}$ step (3): a token count that is not a multiple of
$\mathrm{MAX\_NOISE\_RUN}+1$; a real-token count that is such a multiple but
is not $B+2$ for any reachable bucket $B$; and a padded length prefix
claiming more codepoints than the padded buffer holds. Each positive vector
fixes $(\mathbf{k},\mathsf{sk},A,M)$ and the expected
ciphertext; because $\mathsf{Enc}$ is a pure function of its inputs
(Definition~\ref{def:aead-triple}), no nonce field appears and none is
needed. Conforming implementations must produce byte-identical ciphertext
for every positive vector and reject every negative vector. The generator
\texttt{tests/gen\_kats\_v8.py} regenerates the file and verifies parity on
any port; \texttt{tests/test\_kats.py} and
\texttt{rust/src/kat\_cross\_check.rs} consume it from Python and Rust
respectively.

A second corpus, \texttt{tests/kat/v6\_vectors.json} (37 vectors), is
retained in the repository for the earlier v7 wire format, which takes an
externally supplied nonce and is \emph{not} the scheme analysed here; its
vectors are outside the scope of this paper and none of the results of
Section~\ref{sec:security} are asserted of them.

Both corpora, together with the regression suites across the Python, Rust
and C implementations, also check Proposition~\ref{prop:correctness}
empirically. As stated in Section~\ref{sec:intro}, these are evidence of
cross-implementation agreement only; no security claim in this paper rests
on them.

% ─────────────────────────────────────────────────────────────────────────────
\section{Fuzz Testing}
\label{sec:fuzz}
% ─────────────────────────────────────────────────────────────────────────────

Two fuzz harnesses are provided. The Python harness
(\texttt{tests/fuzz\_atheris.py}, using \texttt{atheris}) targets
\texttt{decrypt\_bytes\_v8} and the fixed-width token decoder, verifying
that no exception other than \texttt{ValueError} escapes the public API.
The Rust harness (\texttt{rust/fuzz/fuzz\_targets/decode\_bytes.rs}, using
\texttt{cargo-fuzz}) targets the Rust \texttt{decrypt\_bytes\_v8} entry
point with multiple key sizes and AAD values. No panics were observed
across $10^6$ corpus entries in either harness.

% ─────────────────────────────────────────────────────────────────────────────
\section{Known Caveats}
\label{sec:caveats}
% ─────────────────────────────────────────────────────────────────────────────

\begin{description}
  \item[Plaintext length cap.] The block padding scheme stores the
        plaintext length as a 2-byte big-endian integer, capping messages
        at 65{,}535 codepoints. Exceeding the cap raises an error
        immediately; there is no silent truncation. Callers requiring
        larger messages must split at the application layer.

  \item[Ciphertext expansion.] Fixed at exactly $20$ tokens, i.e.\ $160$
        bytes, per padded codepoint (Section~\ref{sec:comparison}). The
        expansion is deterministic rather than message-dependent, which is
        what Lemma~\ref{lem:length-det} requires, but it is \emph{not} a
        traffic-analysis countermeasure: by
        Proposition~\ref{prop:expansion-neutral} it contributes zero bits
        of length confidentiality. It is a cost of the arithmetic layer,
        whose retention is argued and bounded in
        Section~\ref{sec:arith-layer-value} and quantified against the same
        construction without that layer in Section~\ref{sec:lean-profile}.

  \item[Padding length-bucket leakage.] Ciphertext length reveals the
        padding bucket, and nothing else about the plaintext
        (Lemma~\ref{lem:length-det}). Under the default $\mathsf{bucket}$
        profile there are $13$ reachable buckets
        $\{16,32,\ldots,65536\}$, so the leakage is at most
        $\log_2 13\approx3.70$ bits per message
        (Corollary~\ref{cor:length-leak}) — it does not reveal the exact
        codepoint count within a bucket, and, because the token ceiling is
        fixed per bucket, it does not vary with $\mathsf{sk}$, the nonce, or
        the associated data, so no averaging attack across repeated
        encryptions of one message under varying AAD can extract more than
        the bucket itself. Deployments requiring strictly less than this
        should select a coarser padding profile, and those requiring zero
        length leakage should select $\mathsf{frame}(F)$
        (Section~\ref{sec:pad-profiles}), which needs no change to the wire
        format and no cooperation from the decryptor; no external
        fixed-frame transport is required.

  \item[Scope of $\mathsf{sk}$ and $\mathbf{k}$ compromise.] $\mathsf{sk}$
        is the sole cryptographic secret in this construction: every HMAC
        domain derivation — the noise schedule, the padding, the
        per-position addends, the keystream, and the tag — is keyed by
        $\mathsf{sk}$ alone (Section~\ref{sec:domains}). Consequently, an
        adversary who compromises $\mathsf{sk}$ but not $\mathbf{k}$ can
        recover any real token's plaintext codepoint directly: the
        per-position addend's raw HMAC output depends only on
        $\mathsf{sk}$, so for a given token the adversary can, for each of
        the $579{,}947$ candidate primes in $\mathcal{P}$, compute
        the corresponding addend and check whether the resulting quotient
        is an exact integer in a valid codepoint range — an
        $\mathcal{O}(|\mathcal{P}|)$ trial-division search that succeeds in
        microseconds. $\mathbf{k}$'s secrecy therefore adds \textbf{no}
        independent confidentiality margin once $\mathsf{sk}$ is known:
        compromise of $\mathsf{sk}$ alone is equivalent to a full break,
        regardless of whether $\mathbf{k}$ also leaked. Operationally,
        $\mathbf{k}$ and $\mathsf{sk}$ must therefore be protected with
        equal rigor, as a single unit — not because splitting them across
        separate storage or trust domains buys compartmentalized
        defense-in-depth (it does not), but because $\mathsf{sk}$
        compromise alone is already fatal. This is consistent with, and
        does not weaken, the security theorems of
        Section~\ref{sec:security}, which assume the entire key material
        is secret from the adversary, exactly as for any symmetric AEAD
        scheme (a leaked AES-256 key is likewise a full break of AES-GCM).

  \item[Message domain and the C port.] The message space $\mathcal{M}$
        excludes the surrogate range $[\mathtt{U+D800},\mathtt{U+DFFF}]$,
        since the synthetic-nonce derivation hashes the UTF-8 encoding of
        $M$ and UTF-8 is undefined there. Rust enforces this by type and
        Python rejects surrogates explicitly. The C port
        (\texttt{C/napqes.c}) instead maps each \emph{byte} of a
        NUL-terminated \texttt{char\,*} to one codepoint
        (Remark~\ref{rem:impl-domain}), so it agrees with the other two
        ports on ASCII input only; on non-ASCII input it produces a
        different, self-consistent ciphertext, and it cannot represent an
        embedded NUL at all. Cross-implementation interoperability is
        therefore claimed for ASCII plaintexts only, which is the domain the
        KAT corpus of Section~\ref{sec:kats} covers. Fixing the C port to be
        codepoint-oriented would change no wire format but is not done here.

  \item[Per-key data limit.] The $q^2/2^{128}$ term in
        Theorems~\ref{thm:ind-cpa} and~\ref{thm:ind-cca} is an ordinary
        birthday bound over the 128-bit synthetic nonce, so NAPQES has a
        data-complexity limit of roughly $2^{64}$ encryptions per key, the
        same order as any 128-bit-nonce AEAD. Synthetic nonce derivation
        removes the \emph{caller-error and DRBG-failure} routes to nonce
        reuse (Definition~\ref{def:synthnonce}); it does not remove the
        birthday bound, and Section~\ref{sec:comparison} should not be read
        as claiming otherwise. We recommend an operational cap of
        $q\le2^{48}$ encryptions per $(\mathbf{k},\mathsf{sk})$ pair, at
        which the term is $2^{-32}$, followed by rekeying. Neither the
        reference implementations nor the wire format enforce this cap;
        it is a deployment obligation.

  \item[Timing side channels.] The security theorems of
        Section~\ref{sec:security} are black-box statements and cover no
        side channel. Only the tag comparison is written to run in constant
        time; the token-emission and token-inversion loops iterate a
        noise-dependent number of times and perform integer division by a
        key element, both of which are data-dependent on real hardware, and
        the Python reference implementation offers no timing guarantees at
        all. Section~\ref{sec:constant-time} states precisely what is and is
        not covered. Deployments in which an adversary can measure
        encryption or decryption time on inputs it influences should treat
        this as an unresolved gap; a TVLA or dudect study of the full
        encode/decode path, rather than of the tag comparison alone, remains
        future work.
\end{description}

% ─────────────────────────────────────────────────────────────────────────────
\section{Conclusion}
\label{sec:conclusion}
% ─────────────────────────────────────────────────────────────────────────────

We have presented NAPQES, an AEAD scheme whose entire security reduces to
the pseudorandomness of HMAC-SHA256. The construction relies solely on
FIPS-approved primitives, and keys two independently sampled secrets — an
arithmetic-layer prime tuple $\mathbf{k}$ and a uniform 256-bit HMAC key
$\mathsf{sk}$ — together with a synthetic, message-bound nonce that removes
nonce selection from the caller and degrades identical-message reuse only
to disclosed plaintext equality. Like every symmetric AEAD scheme it is
unaffected by Shor's algorithm; what distinguishes it from AES-GCM and
ChaCha20-Poly1305 is the absence of a polynomial authenticator and hence of
the nonce-reuse key-recovery mode that follows from one. We prove
IND-CPA-det, INT-CTXT, and IND-CCA-det security under the standard,
uniform-key PRF assumption on HMAC-SHA256, with no key-entropy term and no
simulation gap, and prove that ciphertext length is a deterministic
function of the padded-length bucket alone. Against a length-observing
adversary we prove the strictly stronger LH-IND-CPA-det notion
(Theorem~\ref{thm:lh-ind-cpa}) at the same bound, bound the resulting
leakage at $\log_2 13\approx3.70$ bits per message and zero under the
$\mathsf{frame}$ padding profile (Corollary~\ref{cor:length-leak}), and
exhibit the zero-query adversary that breaks the same notion outright for
AES-GCM and ChaCha20-Poly1305
(Proposition~\ref{prop:lh-separation}). We also show that this property
follows from the padding profile alone: the arithmetic layer contributes
nothing to it (Proposition~\ref{prop:expansion-neutral}), and removing that
layer changes no theorem in this paper
(Proposition~\ref{prop:lean-theorems}). Separately from the proofs, the
Python, Rust and C implementations agree on the cross-implementation KAT
corpora of Section~\ref{sec:kats}, the keystream passes the NIST SP 800-22
Rev.\ 1a suite (Section~\ref{sec:sts}), and the decoders survive the fuzz
harnesses of Section~\ref{sec:fuzz}; these are evidence of implementation
agreement, statistical quality and decoder robustness respectively, and
none of them is evidence for a security claim, all of which rest on the
theorems alone.

Remaining open limitations are documented in Section~\ref{sec:caveats}:
the fixed $160$-byte-per-codepoint ciphertext expansion of the
noise-injection design, which buys no theorem and is retained for the
beyond-model reasons set out in Section~\ref{sec:arith-layer-value}, the
16-bit plaintext length cap, the residual padding length-bucket leak, and
the scope of what $\mathsf{sk}$ and $\mathbf{k}$ compromise each imply.
Future work includes a TVLA side-channel analysis of the constant-time
Rust implementation, a FIPS 140-3 module boundary definition, and a
performance evaluation on embedded ARM Cortex-M4 and RISC-V targets.

% ─────────────────────────────────────────────────────────────────────────────
\bibliographystyle{plain}
\begin{thebibliography}{99}

\bibitem{nist-gcm}
  M. Dworkin,
  \emph{Recommendation for Block Cipher Modes of Operation: Galois/Counter Mode (GCM) and GMAC},
  NIST SP 800-38D, November 2007.

\bibitem{rfc8439}
  Y. Nir, A. Langley,
  \emph{ChaCha20 and Poly1305 for IETF Protocols},
  RFC 8439, June 2018. (Obsoletes RFC 7539.)

\bibitem{rfc5297}
  D. Harkins,
  \emph{Synthetic Initialization Vector (SIV) Authenticated Encryption Using the Advanced Encryption Standard (AES)},
  RFC 5297, October 2008.

\bibitem{rfc8452}
  S. Gueron, A. Langley, Y. Lindell,
  \emph{AES-GCM-SIV: Nonce Misuse-Resistant Authenticated Encryption},
  RFC 8452, April 2019.

\bibitem{ascon2021}
  C. Dobraunig, M. Eichlseder, F. Mendel, M. Schl\"{a}ffer,
  \emph{Ascon v1.2: Lightweight Authenticated Encryption and Hashing},
  Journal of Cryptology, 34(3), 2021.

\bibitem{sp800-232}
  M. S. Turan, K. McKay, D. Chang, L. E. Bassham, J. Kang, N. D. Waller, J. M. Kelsey, D. Hong,
  \emph{Ascon-Based Lightweight Cryptography Standards for Constrained Devices: Authenticated Encryption, Hash, and Extendable Output Functions},
  NIST SP 800-232, August 2025.

\bibitem{sp800-131a}
  E. Barker, A. Roginsky,
  \emph{Transitioning the Use of Cryptographic Algorithms and Key Lengths},
  NIST SP 800-131A Rev.\ 2, March 2019.

\bibitem{grover96}
  L. K. Grover,
  \emph{A fast quantum mechanical algorithm for database search},
  Proceedings of STOC '96, pp. 212--219, 1996.

\bibitem{bbbv1997}
  C. H. Bennett, E. Bernstein, G. Brassard, U. Vazirani,
  \emph{Strengths and Weaknesses of Quantum Computing},
  SIAM Journal on Computing, 26(5), pp. 1510--1523, 1997.

\bibitem{zalka1999}
  C. Zalka,
  \emph{Grover's quantum searching algorithm is optimal},
  Physical Review A, 60(4), pp. 2746--2751, 1999.

\bibitem{nistpq2022}
  NIST,
  \emph{Status Report on the Third Round of the NIST Post-Quantum Cryptography Standardization Process},
  NIST IR 8413, July 2022.

\bibitem{joux2006}
  A. Joux,
  \emph{Authentication failures in NIST version of GCM},
  NIST Comment, 2006.

\bibitem{iwata2012}
  T. Iwata, K. Ohashi, K. Minematsu,
  \emph{Breaking and Repairing GCM Security Proofs},
  CRYPTO 2012, LNCS 7417, pp. 31--49.

\bibitem{rfc5869}
  H. Krawczyk, P. Eronen,
  \emph{HMAC-based Extract-and-Expand Key Derivation Function (HKDF)},
  RFC 5869, May 2010.

\bibitem{nist-drbg}
  E. Barker, J. Kelsey,
  \emph{Recommendation for Random Number Generation Using Deterministic Random Bit Generators},
  NIST SP 800-90A Rev 1, June 2015.

\bibitem{bellare2000ae}
  M. Bellare, C. Namprempre,
  \emph{Authenticated Encryption: Relations among Notions and Analysis of the Generic Composition Paradigm},
  ASIACRYPT 2000, LNCS 1976, pp. 531--545.

\bibitem{rogaway2006}
  P. Rogaway, T. Shrimpton,
  \emph{A Provable-Security Treatment of the Key-Wrap Problem},
  EUROCRYPT 2006, LNCS 4004, pp. 373--390.

\bibitem{sp800-22}
  A. Rukhin, J. Soto, J. Nechvatal, et al.,
  \emph{A Statistical Test Suite for Random and Pseudorandom Number Generators for Cryptographic Applications},
  NIST SP 800-22 Rev 1a, April 2010.

\end{thebibliography}

\end{document}
